\documentclass[11pt,a4paper]{article}
\usepackage[T1]{fontenc}
\usepackage[utf8]{inputenc}
\DeclareUnicodeCharacter{2082}{\textsubscript{2}}
\usepackage{lmodern}
\usepackage[a4paper,margin=24mm,headheight=14pt]{geometry}
\usepackage{amsmath,amssymb,amsthm,mathtools}
\usepackage{booktabs,longtable,array,enumitem}
\usepackage{microtype,xcolor}
\usepackage{xurl}
\usepackage[colorlinks=true,linkcolor=blue!45!black,citecolor=blue!45!black,urlcolor=blue!45!black]{hyperref}
\usepackage{fancyhdr}
\pagestyle{fancy}
\fancyhf{}
\fancyhead[L]{\small ZETA\_TWO: schemes, motives, and the two RH inputs}
\fancyhead[R]{\small 21 November 2026}
\fancyfoot[C]{\thepage}
\setlength{\emergencystretch}{2em}
\setlength{\parskip}{0.25em}
\setlist{itemsep=0.35em,topsep=0.4em}
\newtheorem{theorem}{Theorem}[section]
\newtheorem{proposition}[theorem]{Proposition}
\newtheorem{corollary}[theorem]{Corollary}
\newtheorem{lemma}[theorem]{Lemma}
\theoremstyle{definition}
\newtheorem{definition}[theorem]{Definition}
\newcommand{\Gal}{\operatorname{Gal}}
\newcommand{\Spec}{\operatorname{Spec}}
\newcommand{\Tr}{\operatorname{Tr}}
\newcommand{\Frob}{\operatorname{Frob}}
\title{\textbf{ZETA\_TWO}\\[0.5em]
Riemann hypotheses for schemes and motives\\
from the stipulated Alp\"oge theorems\\
and the conditional Masullo claim}
\author{Research synthesis and proof checks: GPT-6 Astra Pro\\
\small Run ZETA\_TWO-20261121; reasoning effort not exposed}
\date{21 November 2026 (Europe/Berlin)}
\begin{document}
\maketitle
\begin{abstract}
We determine exact sufficient classes of arithmetic schemes and motives to
which the two specified zero theorems transfer. Complete abelian
Artin--Tate point counts reduce to primitive Dirichlet functions, using
Alp\"oge as stipulated. Conditional on Masullo's claim for every unitary
Hecke character over every number field, Brauer induction gives divisor RH
for all Artin factors, and the same-weight Taniyama-group factorization
gives it for every pure potentially CM motive in the absolute-Hodge category.
Neither quotient argument proves general holomorphy.

For an arbitrary finite-type scheme over $\mathbf Z$, a finite
Grothendieck-group reduction isolates pure geometric $L$-functions and
explicit exceptional local factors. A non-CM elliptic curve already
supplies an exact remaining global obstruction: under the stipulated
Dirichlet RH its off-line $L$-zeros cannot cancel from the scheme zeta
divisor. The invariant gate is non-CM geometry up to Tate twist, rather
than positive weight alone. We audit every supplied flag and intake
occurrence, extract effective class-number, prime-witness, and CM-height
bounds, and distinguish them from unconditional special-value arithmetic.
Exact computation certifies the $23$-cyclotomic class number. A mixed Tate
counterexample shows that identical Euler products need not recover
extension classes. Original archive scans are used only where actually
retrieved; unchecked intermediate claims remain visibly separate.
\end{abstract}

\paragraph{Scope of the certification.}
Every new proof in Part I has been read in full and its check recorded in
\path{CHECK_EVENTS.tsv}. Independent model readers checked the principal
transfer proofs and the arithmetic derivations. These are mathematical
proof reviews, not Lean certificates. The original proofs of cited
literature results are background; checking their statements and
hypotheses is recorded separately. No execution of the full
Alp\"oge Lean archive is claimed. Targeted official source files were
retrieved and inspected separately from a kernel build. Accepting its five endpoint theorems is
the task's standing instruction. Masullo's global claim remains an
explicit condition throughout.

\paragraph{How to read the accompanying records.}
\path{QUESTIONS.tsv} gives the result or current source gate for every
treated question. The 575-row flag audit and 587-row intake inventory are
exhaustive audits of supplied metadata, not 1,162 newly certified proofs.
The 85 scoped conditional records distinguish an original row from an
additional estimate one might attach to it. \path{IMPLICATIONS.md} is the
programme-facing synthesis. \path{SOURCE_CORRECTIONS.tsv} corrects intake
prose; \path{TRANSCRIPTION_NOTES.tsv} is reserved for original manuscript
readings. Neither a missing original source nor a new RH consequence
silently reverses an intake's existing ``refuted'' verdict.

\begingroup\small\setlength{\parskip}{0pt}
\tableofcontents
\endgroup
\clearpage
\part{Checked findings and complete transfer proofs}


\section{The main answer, with the two inputs separated}
\label{sec:main-answer}

The results below are transfer theorems, relative to the inputs specified in
the question.  They do not assert a new verification of either proposed
proof of RH.  Throughout, \textbf{A} means the five final declarations in
\path{ALPOGE_Main.lean}, with their meanings in
\path{ALPOGE_Statements.lean}, cited as theorems exactly as requested.
In particular A supplies RH and GRH for Dirichlet characters over
$\mathbf Q$.  \textbf{M} means GRH for the Hecke $L$-functions of
\emph{all continuous unitary} idele-class characters over every number
field, as an explicit assumption on Masullo's draft claim.  The distinction
between its intended claim and its unfinished numbered statement is
documented below.  In the dependency tables, ``A'' means no further
unproved input beyond the stipulated Alp\"oge theorems; ``M'' means
conditional on M; ``U'' means previously unconditional; and ``B'' means
not supplied by A or M.  These labels are not claims of logical
independence of mathematical conjectures.

\begin{center}\small
\begin{tabular}{>{\raggedright\arraybackslash}p{.27\linewidth}>{\raggedright\arraybackslash}p{.28\linewidth}>{\raggedright\arraybackslash}p{.35\linewidth}}
\toprule
Object & Conclusion from A & Additional conclusion under M \\
\midrule
Schemes with exact abelian Artin--Tate counts at every prime
& Complete zeta factorization; noninteger divisor support on
$\Re s=i+1/2$ & Same conclusion \\
Schemes with general Artin--Tate counts
& Abelian part only & Complete factorization by shifted Artin functions;
critical-strip divisor RH for each weight \\
$\operatorname{Spec}\mathcal O_K$
& RH when $K/\mathbf Q$ is abelian & RH for every number field $K$ \\
Artin motives and Tate twists
& Abelian Artin part & All Artin motives; meromorphic divisor RH,
without general Artin holomorphy \\
Pure CM motives in the absolute-Hodge category
& The abelian Artin--Tate subcategory & All of them, using same-weight
Hecke factorization; line $\Re s=(w+1)/2$ \\
Arbitrary finite-type schemes over $\mathbf Z$
& Reduction to a finite list of pure geometric $L$-functions plus explicit
local factors & The Artin--CM factors leave the residual list;
non-CM geometric RH remains \\
\bottomrule
\end{tabular}
\end{center}

No assertion in this table identifies an arbitrary integral model's
zeta function with its generic fiber's canonical finite motivic Euler product.
The difference is an explicit finite product, and it matters for zeros.
Nor does a point-count identity imply the existence of algebraic
projectors giving a decomposition of a Chow motive.

\subsection{Conventions: sheaves, weights, and the meaning of RH}
We use geometric Frobenius on cohomology, so $\mathbf Q_\ell(-1)$ has
Frobenius eigenvalue $p$ and weight $2$.  For a pure motive $N$ of weight
$w$, the Tate twist satisfies
\[
 \operatorname{wt}N(n)=w-2n,\qquad L(N(n),s)=L(N,s+n).
\]
The unitary variable is $z=s-w/2$ and its critical line is $\Re z=1/2$.
Thus the motivic line is $\Re s=(w+1)/2$.  We mean by
\emph{divisor RH} that the meromorphic function has neither a zero nor
a pole off this line inside $w/2<\Re s<w/2+1$.
For functions known to be holomorphic there this is ordinary RH.
Zeros forced by archimedean gamma factors, endpoint poles, and zeros
of removed finite Euler factors must be specified separately.

Deligne's 1979 paper \emph{Valeurs de fonctions $L$ et p\'eriodes
d'int\'egrales}, \S\S1.1--1.3 and (3.1.2), supplies the motivic Euler
factor and twist conventions.  Its Conjectures 1.8 and 2.8 concern
critical special values and periods; we do not attribute to either of
those numbered conjectures the RH question in the prior run's Q02.
That question is retained as the programme's question, with the above
precise analytic meaning.

% Prepared by /root/masullo_cm, GPT-6 Astra Pro; effort not exposed.
% 21 November 2026, Europe/Berlin.
% Bibliography/reference material for the root to integrate; not standalone.

\paragraph{Classical analytic background for the Hecke factors.}
The required continuation and functional equation are classical and do not
depend on the Masullo draft.  A checked precise reference is Poonen,
\emph{Tate's thesis}, \S5.11, Theorem~5.22, pp.~38--39, with its
completed-factor convention in Definition~5.20.  Local factors are given
in \S4.7, p.~19.  For a finite-order Hecke character, the real factors are
$\Gamma_{\mathbf R}(s+\epsilon_v)$, $\epsilon_v\in\{0,1\}$, and the complex
factors are $\Gamma_{\mathbf C}(s)$: use Corollary~4.15 and Remark~4.16,
p.~18.  Their poles lie at real nonpositive integers and they have no
zeros (\S1.1, p.~2).
Nonvanishing on $\Re s=1$ follows, for example, by taking
$n=m=1$, $\pi=\chi$, and $\pi'=1$ in Shahidi,
\emph{On nonvanishing of $L$-functions}, the unnumbered theorem on p.~462.
The finite-order principal character has its usual pole at $1$.
The finitely many local factors excluded in Shahidi's notation have no
zeros or poles on $\Re s=1$, so the result applies to the complete finite
Hecke product used here.


\section{Artin--Tate point counts and their complete Euler products}
\label{sec:artin-tate}

\subsection{The exact class and its geometric interpretation}
Let $R_{\rm Art}(\mathbf Q)$ be the Grothendieck group of finite-image
characteristic-zero representations of $G_{\mathbf Q}$, with coefficients
embedded in $\mathbf C$.  For a virtual representation $V$ and a prime
$p$, write
\[
 a_{p,r}(V)=\operatorname{Tr}(F_p^r\mid V^{I_p}).
\]
Here invariants and trace are extended additively; exactness of invariants
is valid because the relevant inertia action factors through a finite
group in characteristic zero.  Use the reciprocity convention in which
the associated one-dimensional Artin factor equals the indicated
primitive Dirichlet factor.  Inverting all reciprocity characters would
change none of the assertions about their class or zero locations.

\begin{definition}[AT point counts]
A finite-type scheme $X/\mathbf Z$ has \emph{exact Artin--Tate point
counts} if there are finitely many $V_i\in R_{\rm Art}(\mathbf Q)$,
$i\geq0$, such that, for every prime $p$ and integer $r\geq1$,
\begin{equation}\label{eq:ATcounts}
 \#X(\mathbf F_{p^r})=\sum_i p^{ir}a_{p,r}(V_i).
\end{equation}
The class is \emph{abelian} if every constituent of every $V_i$ factors
through a finite abelian quotient of $G_{\mathbf Q}$.
The analogous definition \emph{away from $S$} imposes
\eqref{eq:ATcounts} only for $p\notin S$.
\end{definition}

The coefficients are allowed to be virtual because compactly supported
cohomology has alternating signs.  This is an exact sufficient class,
not a claimed classification of all schemes whose zeta function could
satisfy RH.  A useful sheaf-theoretic sufficient condition is
\begin{equation}\label{eq:sheafclass}
 [Rf_!\mathbf Q_\ell]=\sum_i[\mathcal V_i(-i)]
 \quad\hbox{in }K_0(D^b_c(U,\mathbf Q_\ell)),
 \qquad U=\operatorname{Spec}\mathbf Z[1/S\ell],
\end{equation}
where the $\mathcal V_i$ have finite monodromy.  Proper-support base
change and the Frobenius trace formula give \eqref{eq:ATcounts} on $U$.
This is the natural language of SGA 4, XVII, \S\S5.1--5.3, and SGA 5,
XV, \S3, n\textsuperscript{o}~2, unnumbered theorem, equation (1);
equivalently Grothendieck's Bourbaki
expos\'e 279, Theorem 5.1.  At an exceptional prime one must compare the
actual fiber with the inertia-invariant formula; it is not determined
by \eqref{eq:sheafclass}.  Since a finite-type scheme admits a finite
stratification by separated affine pieces, the same count statements
also cover nonseparated schemes by additivity.  Nilpotents do not alter
points over finite fields or the closed-point Euler product.

For example, a finite disjoint stratification into pieces
$\mathbf A^{i}_{\mathcal O_K}$, with $K$ abelian, has abelian AT counts.
More generally the same holds with arbitrary $K$ for AT counts.  These
examples include the usual split cellular schemes.  Split tori have
counts $(p^r-1)^d$, so their virtual signs are essential.  Products of
strata with split tori, and finite disjoint unions, are covered in the
same way.  This assertion concerns the indicated schemes with their
actual integral models; an arbitrary model with the same generic fiber
need only satisfy the away-from-$S$ version.

\begin{theorem}[ZT01: exact Euler-factor transfer]\label{ZT01}
If $X$ has \eqref{eq:ATcounts}, then
\begin{equation}\label{eq:ATfactor}
 \zeta_X(s)=\prod_i L(V_i,s-i)
\end{equation}
first in a right half-plane, and then meromorphically on $\mathbf C$.
This factorization is unconditional.  In the abelian case it is a
finite product
\[
 \zeta_X(s)=\prod_{i,\chi}L(s-i,\chi)^{m_{i,\chi}}
\]
of primitive Dirichlet functions, with integer exponents.
Under A every zero or pole of this exact abelian product is either a
real integer or lies on one of the lines $\Re s=i+1/2$.
\end{theorem}

\begin{proof}
For a finite-dimensional vector space with endomorphism $T$, the
identity
\begin{equation}\label{eq:logdet}
 -\log\det(1-Tt)=\sum_{r\geq1}\operatorname{Tr}(T^r)\frac{t^r}{r}
\end{equation}
holds as a formal power series: over an algebraic closure both sides
are the sum of $-\log(1-\alpha t)$ over the eigenvalues, counted with
algebraic multiplicity.  Subtraction extends the identity to virtual
representations.  The elementary closed-point identity over a finite
field gives
\[
 Z(X_p,t)=\exp\left(\sum_{r\geq1}
          \#X(\mathbf F_{p^r})\frac{t^r}{r}\right).
\]
Indeed a closed point of degree $d$ contributes $d$ points over
$\mathbf F_{p^r}$ precisely when $d\mid r$; summing its contribution
gives $\sum_{k\geq1}t^{dk}/k=-\log(1-t^d)$.
Substitute \eqref{eq:ATcounts} into this identity and apply
\eqref{eq:logdet}.  The result at each prime is
\[
 Z(X_p,p^{-s})
   =\prod_i\det(1-p^{i-s}F_p\mid V_i^{I_p})^{-1}.
\]
If the expression is empty the identity is immediate. Otherwise let $d$ be the largest $i$ occurring and let $C$ be the sum of the
dimensions with absolute multiplicities of all the representations.
Finite-image eigenvalues have absolute value one.  Thus the absolute
value of the summand of the logarithmic series is bounded by
$Cp^{-r(\Re s-d)}/r$.  Its sum over primes and $r$ converges when
$\Re s>d+1$.  Consequently products and sums can be interchanged there,
giving \eqref{eq:ATfactor}.  Meromorphic continuation follows from
Brauer induction and classical Hecke continuation for Artin $L$-functions.

In the abelian case, split each $V_i$ over $\mathbf C$ into
one-dimensional finite-order characters.  Global abelian reciprocity
over $\mathbf Q$ identifies their complete Artin Euler products with
primitive Dirichlet $L$-functions.  At a ramified prime a character has
zero-dimensional inertia invariants; its Euler factor is $1$, exactly
as the primitive Dirichlet character has value $0$ there.  Thus the
identity includes all bad primes and no factor has been silently lost.

A primitive nonprincipal Dirichlet $L$-function is entire, with
gamma-forced trivial zeros at real nonpositive integers; its remaining
zeros lie in the critical strip.  The principal primitive factor is
$\zeta(s)$, with its pole at $1$.  A places each nontrivial zero on
$\Re s=1/2$.  Shifting by $i$ puts the nontrivial support on
$\Re s=i+1/2$ and the other divisors at integers.  Finally, at any point,
the order of a finite meromorphic product is the integer sum of the
orders of its factors.  Cancellation can remove support but cannot
create support at a point where every factor is invertible.  This
proves the last assertion, including negative exponents.
\end{proof}

\begin{theorem}[ZT02: general Artin transfer under M]\label{ZT02}
Assume M.  For every virtual Artin representation $V$ over $\mathbf Q$,
every zero or pole of $L(V,s)$ in $0<\Re s<1$ lies on $\Re s=1/2$.
For $V(-i)$ the corresponding strip is $i<\Re s<i+1$ and the line is
$\Re s=i+1/2$.  In \eqref{eq:ATfactor}, every noninteger zero or pole
is on one of these lines.  No Artin holomorphy assertion is required
or obtained.
\end{theorem}

\begin{proof}
Choose a finite Galois extension $L/\mathbf Q$ through which $V$ factors
and put $G=\operatorname{Gal}(L/\mathbf Q)$.  Brauer induction gives
an equality of virtual characters
\[
 [V]=\sum_a n_a[\operatorname{Ind}_{H_a}^G\theta_a],
 \qquad n_a\in\mathbf Z,
\]
where $H_a\subset G$ and each $\theta_a$ is a finite-order
one-dimensional character of $H_a$.  Put $F_a=L^{H_a}$.
Artin induction formalism and global reciprocity over $F_a$ give the
complete Euler-product identity
\[
 L(V,s)=\prod_a L_{F_a}(\theta_a,s)^{n_a}.
\]
These two classical inputs may be found together, for example, in
Anderson \cite[\S3.8, Proposition 3.8.1; \S5.7, Brauer Induction
Lemma]{Anderson1986}, specialized to a finite group, and the
rank-one reciprocity interpretation in \cite[\S7.2]{FarguesCM}.
Each finite-order Hecke factor is holomorphic in the open critical
strip, and M makes it nonzero off the central line.  The valuation
identity for products proves the strip assertion.  The classical
functional equation and gamma factors for finite-order Hecke
characters put their other divisors at real integers, including the
possible pole at $1$; the gamma factors have integral shifts.
Taking an integer quotient cannot introduce support elsewhere.
Tate translation then proves the remaining assertions and combines
with ZT01.  A negative exponent may leave a pole at a common critical
zero.  Zero location alone gives no inequality on that signed order,
which is why the argument does not prove holomorphy.
\end{proof}

\begin{corollary}[ZT03: all number fields, and the abelian specialization]
\label{ZT03}
For every number field $K$, $X=\operatorname{Spec}\mathcal O_K$ has
exact AT counts, and its factor in ZT01 is $\zeta_K(s)$.  Hence M
implies RH for every Dedekind zeta function.  If $K/\mathbf Q$ is
abelian, A already gives the conclusion and the primitive Dirichlet
factorization of prior-run F2 is valid at ramified primes as well.
\end{corollary}

\begin{proof}
Let $\Omega=\operatorname{Hom}_{\mathbf Q}(K,\overline{\mathbf Q})$
and $V=\mathbf C[\Omega]$.  For a decomposition group at $p$,
its orbits on $\Omega$ correspond to primes $\mathfrak p\mid p$ in
$K$.  In the orbit for $\mathfrak p$, inertia has $f_{\mathfrak p}$
orbits, where $f_{\mathfrak p}$ is the residue degree, and Frobenius
permutes these in a cycle of length $f_{\mathfrak p}$.
The inertia invariants of the permutation representation have the
orbit sums as a basis.  Their Frobenius trace is therefore
\[
 a_{p,r}(V)=\sum_{\mathfrak p\mid p}
 f_{\mathfrak p}\,\boldsymbol 1_{f_{\mathfrak p}\mid r}
          =\#\operatorname{Spec}\mathcal O_K(\mathbf F_{p^r}).
\]
This argument includes ramification: nilpotents in the special fiber
do not add maps to a field.  The determinant for a cycle of length
$f$ is $1-t^f$, so $L(V,s)=\prod_{\mathfrak p}
(1-(N\mathfrak p)^{-s})^{-1}=\zeta_K(s)$.  Apply ZT02.
When $K/\mathbf Q$ is abelian, $V$ is the regular representation of
its abelian Galois group and splits as the sum of all its characters.
ZT01 applies with each multiplicity equal to one.  Dedekind zeta is
already known to be holomorphic in its open critical strip, so its
divisor RH is the usual zero-only statement.
\end{proof}

\subsection{Artin--Tate motives and the role of the Grothendieck group}
ZT01 and ZT02 apply without a point-count hypothesis to a motive whose
realization is a direct sum of $V_i(-i)$: for motives, $i$ may be any integer; the restriction $i\geq0$ belongs only to the scheme count class. Its $L$-function is the
corresponding product by definition of the local factors.  For an
honest pure abelian Artin--Tate motive of weight $2i$, the character
multiplicities are nonnegative, so holomorphy in the critical strip
is classical.  General Artin constituents give only the meromorphic
statement under M.  In a compact-support Euler characteristic the
multiplicities may have either sign, hence zeros and poles are equally
natural.  For mixed objects this is a statement weight by weight,
not a single RH line for the whole object.

The count hypothesis concerns the virtual realization.  If equalities
of traces are all that is given, the Euler product still follows by
ZT01; an algebraic motivic decomposition is not necessary.  In
particular, no standard conjecture, Tate conjecture, or identification
of numerical and homological equivalence is hidden in that proof.

\section{Finite Euler factors: a necessary correction to a global slogan}
\label{sec:finite-factors}

\begin{theorem}[ZT04: the exact correction outside a good open]
\label{ZT04}
Suppose \eqref{eq:ATcounts} holds only for $p\notin S$, with $S$ finite.
Set
\[
 L_p(V_i,s-i)=\det(1-p^{i-s}F_p\mid V_i^{I_p})^{-1}.
\]
Then
\begin{equation}\label{eq:Rfactor}
 \zeta_X(s)=R_{X,S}(s)\prod_i L(V_i,s-i),\qquad
 R_{X,S}(s)=\prod_{p\in S}
  \frac{Z(X_p,p^{-s})}{\prod_i L_p(V_i,s-i)}.
\end{equation}
Each factor of $R_{X,S}$ is rational in $p^{-s}$.
Its divisor can include entire vertical arithmetic progressions; it
is not in general confined to real integers.
\end{theorem}

\begin{proof}
Repeat the absolutely convergent calculation of ZT01 over $p\notin S$.
Restore the actual local zeta factors for $p\in S$ on the left and
the canonical Artin factors on the right.  Their ratio is exactly
\eqref{eq:Rfactor}.  Rationality of $Z(X_p,t)$ is the finite-field
trace formula (SGA 5, XV, \S3, n\textsuperscript{o}~2, unnumbered
theorem, equation (1)), without an RH
assumption.  If a nonconstant factor is $1-\alpha p^{-s}$ with
$\alpha\ne0$, its zeros satisfy
\[
 \Re s=\frac{\log|\alpha|}{\log p},\qquad
 \Im s=\frac{\arg\alpha+2\pi k}{\log p},\quad k\in\mathbf Z.
\]
The same support description applies to a pole if the exponent is
negative.  This proves the statement and gives an explicit procedure
for locating every exceptional local divisor once the finitely many
local polynomials are specified.
\end{proof}

\begin{proposition}[ZT05: counterexample to ignoring omitted primes]
\label{ZT05}
The assertion that every noninteger zero of an arbitrary integral
model with Tate generic fiber lies on a global RH line is false.
For $X=\operatorname{Spec}\mathbf Z[1/p]$,
\[
 \zeta_X(s)=(1-p^{-s})\zeta(s)
\]
has zeros at $s=2\pi i k/\log p$ for every $k\ne0$.
\end{proposition}

\begin{proof}
The closed points of $X$ omit exactly the prime $p$, giving the
displayed Euler product.  At the indicated nonzero purely imaginary
points, $1-p^{-s}=0$, whereas $\zeta(s)$ has no pole (its only pole
is at $1$).  Thus the product has a zero there, whether or not the
second factor also vanishes.  These points are not real integers
and have real part zero, rather than the Tate global line $1/2$.
\end{proof}

The original F1 in \path{BRIDGES_20261120.tex}, \S3, has an
\emph{all-prime primitive-character} hypothesis and is not refuted
by this example.  The example refutes an extension of its conclusion
to an arbitrary good-open identity.  It also explains why using an
imprimitive principal character without recording its missing Euler
factor would be an error.  For general finite fibers, Weil II,
Theorem 3.3.1, puts their cohomological eigenvalues at integral weights,
so their local divisor lines are $\Re s=v/2$.  These are local
weight lines; global motivic RH concerns $(w+1)/2$.  The shift of
one is essential.

\section{A finite reduction of Q01 to pure geometric factors}
\label{sec:general-gate}

\begin{theorem}[ZT06: a precise scheme-by-scheme gate]\label{ZT06}
For every finite-type $X/\mathbf Z$ there exist finitely many smooth
projective varieties $Y_a/\mathbf Q$, integers $n_a$, and a finite
set $S$ such that, on a right half-plane,
\begin{equation}\label{eq:general-reduction}
 \zeta_X(s)=\left(\prod_{p\in S}Z(X_p,p^{-s})\right)
 \prod_{a,j} L^S(H^j(Y_a),s)^{n_a(-1)^j}.
\end{equation}
Every $H^j(Y_a)$ is pure of weight $j$ at good primes.  If the finitely
many $L$-functions in this formula have meromorphic continuations,
the identity continues meromorphically.  Suppose in addition that
each canonical completed factor has its standard functional equation,
its trivial and boundary divisors are specified, and its remaining
divisor in the open critical strip satisfies divisor RH.  Then all
divisors of $\zeta_X$ are supported on the corresponding global RH
lines together with these specified divisors and the finite local
divisors.  A strip statement alone is not being used to control a
function outside its strip.
Under M every known CM factor in this finite list satisfies its RH
condition; the remaining conditions concern only the residual
non-CM geometric factors and any analytic continuation not already
known for them.
\end{theorem}

\begin{proof}
We may replace $X$ by its reduction and use a finite affine
stratification, since zeta functions multiply under a closed--open
decomposition.  Over $\mathbf Q$, the characteristic-zero
Grothendieck group of varieties is generated by smooth projective
varieties (Bittner, Theorem 3.1, projective refinement).  Thus
\[
 [X_{\mathbf Q}]=\sum_a n_a[Y_a]\quad\hbox{in }K_0(\operatorname{Var}_{\mathbf Q}).
\]
Only finitely many equations, morphisms, and scissor relations occur
in an expression witnessing this equality.  Choose integral models
for them and invert a common nonzero integer to clear their
coefficients and all equations witnessing inverses or isomorphisms.
The models of the $Y_a$ may be made smooth projective by increasing
$S$, including a chosen auxiliary $\ell$: a projective embedding
and its equations spread, and the
nonsmooth locus has empty generic fiber, so properness of that
locus over the base leaves at most finitely many excluded primes.
Likewise the finitely many scissor identities and isomorphisms
spread.  This is the usual finite-presentation limit argument of
EGA IV$_3$, (8.8.2), with the standard openness of the smooth locus.
Consequently the displayed equality holds in
$K_0(\operatorname{Var}_{\mathbf F_p})$ for every $p\notin S$,
and hence gives equality of point counts over every
$\mathbf F_{p^r}$.

For each smooth projective model, smooth proper base change and
the finite-field trace formula give
\[
 \#Y_a(\mathbf F_{p^r})
 =\sum_j(-1)^j\operatorname{Tr}(F_p^r\mid H^j(Y_{a,\overline{\mathbf Q}},
                                             \mathbf Q_\ell)).
\]
Apply \eqref{eq:logdet} at every good prime.  Deligne's Weil I,
Theorem 1.6, supplies $|\alpha|=p^{j/2}$ for these eigenvalues,
so all the logarithmic Euler series converge absolutely sufficiently
far to the right.  This gives \eqref{eq:general-reduction}.
The finitely many exceptional fibers are already recorded exactly.
Uniqueness of meromorphic continuation proves the continuation
assertion whenever its factors have such continuations.  The order
of the product at a point is the signed sum of the factor orders,
so support is contained in the union of the factor supports.
Absolute convergence gives nonvanishing to the right of the
critical strip; the assumed functional equations and the specified
trivial and boundary divisors provide the outside-strip part of
that support statement.
CM01 below settles the CM factors under M.  Removing any factor
already evaluated this way leaves precisely the asserted residual
product and its still-required analytic properties.
\end{proof}

The theorem is a reduction to a finite explicit type of obstruction
for each $X$; it is not a claim that arbitrary motivic analytic
continuation has been proved.  It does not need an unproved
algebraic K\"unneth projector for each $H^j$: these factors can be
read as geometric compatible systems.  When a chosen category of
motives supplies the projectors, they can be denoted $h^j(Y_a)$.
Nor is a converse from a product to its factors claimed: divisors
of distinct factors can cancel.

\begin{proposition}[ZT07: the elliptic-curve obstruction is already present]
\label{ZT07}
Let $E/\mathbf Q$ be an elliptic curve and let $\mathcal E$ be its
smooth proper model over $U=\operatorname{Spec}\mathbf Z[1/S]$.
Then
\[
 \zeta_{\mathcal E}(s)=\frac{\zeta^S(s)\zeta^S(s-1)}{L^S(E,s)}.
\]
For a CM curve, M settles the global nontrivial divisor lines of
this product.  For a non-CM curve the remaining global zero-location
assertion is RH for $L(E,s)$ on $\Re s=1$, even though its analytic
continuation and functional equation are already known.
\end{proposition}

\begin{proof}
At a good prime, $H^0$ has eigenvalue $1$, $H^2$ has eigenvalue $p$,
and $H^1$ has the two roots $\alpha_p,\beta_p$ of
$T^2-a_pT+p$.  The trace formula and \eqref{eq:logdet} yield
\[
 Z(E_p,p^{-s})=
 \frac{(1-\alpha_pp^{-s})(1-\beta_pp^{-s})}
      {(1-p^{-s})(1-p^{1-s})}.
\]
Multiplication over good primes gives the formula.  CM01 handles
$H^1$ in the CM case.  In the non-CM case neither the finite-image
Artin theorem nor the potentially CM Hecke factorization applies
to $H^1(E)$; the required remaining assertion is exactly the
indicated RH, not a missing finite-field Weil bound.

Under A the obstruction in ZT07 is an exact obstruction in the open strip
$1/2<\Re s<3/2$, away from $\Re s=1$.
Indeed $\zeta^S(s)=\zeta(s)\prod_{p\in S}(1-p^{-s})$
has no zero or pole there off the line $\Re s=1$:
the nontrivial global zeros are on the boundary $\Re s=1/2$,
the deleted Euler factors have zeros on $\Re s=0$, and its only pole is
the central point $1$. Similarly $\zeta^S(s-1)$ has global nontrivial
zeros on the other boundary $\Re s=3/2$, deleted-factor zeros on
$\Re s=1$, and its pole at $2$ outside the strip.
The trivial zeros of both zeta factors are outside it.
Furthermore
\[
 L^S(E,s)=L(E,s)\prod_{p\in S}P_p(E,p^{-s}).
\]
At a removed good prime the roots of $P_p$ in the $s$-variable lie on
$\Re s=1/2$; at a multiplicative-reduction prime the polynomial is
$1-a_p p^{-s}$ with $a_p=\pm1$ and its roots lie on $\Re s=0$;
at an additive-reduction prime the polynomial is $1$.
(These are the local factors in Silverman, \emph{The Arithmetic of
Elliptic Curves}, Exercise VIII.8.19 and part~(b).)
Hence every such local correction is invertible in the open strip.
For each $s_0$ in that strip off $\Re s=1$, therefore,
\[
 \operatorname{ord}_{s=s_0}\zeta_{\mathcal E}(s)
       =-\operatorname{ord}_{s=s_0}L(E,s).
\]
Thus a hypothetical off-line zero of the elliptic $L$-function cannot
be concealed by cancellation in the scheme zeta function. This is
stronger than merely observing the elliptic factor in the product.

For the bad-prime local polynomials used here, see Silverman,
\emph{The Arithmetic of Elliptic Curves}, Exercise VIII.8.19.
\end{proof}

Serre's open-image theorem, \S4.4, Th\'eor\`eme~3, p.~299 of \emph{Propri\'et\'es galoisiennes des points d'ordre fini des courbes elliptiques}, makes the non-CM obstruction precise: the connected Zariski monodromy is $\operatorname{GL}_2$, and remains nonabelian after finite extension.

For a holomorphic cuspidal eigenform of weight $k\geq2$, the
associated pure geometric system has weight $k-1$, so its desired
line in motivic normalization is $\Re s=k/2$.  Forms with CM are
included in CM01.  Modularity and automorphy give analytic structure,
but the RH for a non-CM form is an additional assertion.  Weight-one
Artin forms fall in the Artin class and should not be included among
the positive-weight non-CM examples.  More intrinsically the gate
is the geometric part outside the potentially abelian/CM category,
up to Tate twists; it is not the numerical condition $w>0$.

\section{The lambda-ring description and the limit of the local analogy}
\label{sec:lambda}

\begin{proposition}[ZT08: the determinant is the operative lambda invariant]
\label{ZT08}
Let $V$ be a finite-dimensional characteristic-zero vector space and
$F$ an invertible operator.  View the pair $(V,F)$ as a representation
of the infinite cyclic group, and take lambda and Adams operations
in its representation ring.  The local factor is determined by
\[
 \det(1-tF\mid V)=\sum_{j\geq0}(-t)^j
                 \operatorname{Tr}(F\mid\wedge^jV),
 \qquad
 \operatorname{Tr}(F\mid\psi^r V)=\operatorname{Tr}(F^r\mid V),
\]
where $\psi^r$ is the Adams operation on the representation ring.
Thus ZT01 is an identity of determinants arising from a virtual
Artin--Tate realization; it depends on traces and Tate weights,
and not on extension classes or chosen integral lattices.
\end{proposition}

\begin{proof}
Over an algebraic closure, let the eigenvalues of $F$ be
$\alpha_1,\ldots,\alpha_d$.  The trace on $\wedge^jV$ is the
$j$th elementary symmetric polynomial in these eigenvalues, so
the first formula is the expansion of $\prod_a(1-\alpha_at)$.
By the defining Newton identities for Adams operations, the
character of $\psi^rV$ is the $r$th power sum $\sum_a\alpha_a^r$,
giving the second formula.  Additivity extends the logarithmic
formula to virtual representations.  A short exact sequence puts
$F$ in block triangular form, whose determinant is the product of
the two diagonal determinants.  Changing an invariant lattice
does not change the rational operator.  This proves the final
assertion.  At a ramified prime one first takes $V^{I_p}$; one
must not replace $\wedge^j(V^{I_p})$ by $(\wedge^jV)^{I_p}$
without a further argument, since invariants are not generally
monoidal.
\end{proof}

This is the part of Grothendieck's lambda formalism actually used.
Crystalline Frobenius can compute the same good-reduction local
polynomial when the relevant comparison theorem applies.  It does
not turn a bound on local eigenvalues into a theorem about zeros
of the infinite product.  Divided powers and integral theta
derivatives concern additional integral structure; neither ZT01
nor ZT08 supplies an integral crystal, a divided-power stratification,
or the theta divisibility bound claimed in the note's step 2.
Those require their own hypotheses and proofs.  This keeps the
global transfer theorems separate from the other parallel prompts
on crystals, multiplicity, and the standard conjectures.


% Research contribution by /root/masullo_cm, 2026-11-21 Europe/Berlin.
% Bibliography keys: Anderson1986, FarguesCM, Milne1990.
\section{CM motives: the same-weight Hecke reduction}

\paragraph{Conventions and the hypothesis actually used.}
Let $\mathsf{PCM}(\mathbf Q)$ denote the category of motives for absolute
Hodge cycles generated by Artin motives, Tate motives, and the cohomology of
abelian varieties which acquire complex multiplication over a number field.
It is closed under direct sums, summands, tensor products, and duals.
This is Anderson's category in \cite[\S5.3]{Anderson1986}; the term CM motive
below means an object of this category. For a pure object $M$ of weight $w$,
geometric Frobenius eigenvalues at good primes have absolute value $p^{w/2}$.
Our twist convention is
\[
  \operatorname{wt}M(n)=w-2n,\qquad L(M(n),s)=L(M,s+n).
\]
In particular $V(-i)$ for an Artin representation $V$ has weight $2i$ and
$L(V(-i),s)=L(V,s-i)$.

Write \textbf{(M)} for the following assumption, abstracted from the draft's
intended claim: for every number field $F$ and every continuous unitary
idele-class character $\chi$, every zero of its completed Hecke $L$-function
in $0<\Re z<1$ has $\Re z=1/2$. This includes characters of infinite order.
The draft uses $\widehat C_F$ for the unitary dual (PDF p.~7), and its
\S1.2 (PDF p.~12) takes an arbitrary $\chi\in\widehat C_F$.
Thus a hypothesis restricted to finite-order Hecke characters would not
suffice for the CM conclusion below.

\paragraph{Finding CM01 (conditional divisor RH for every pure CM motive).}
Assume \textbf{(M)}. Let $M\in\mathsf{PCM}(\mathbf Q)$ be pure of weight $w$.
Its $L$-function is meromorphic on $\mathbf C$. Every zero or pole of
$L(M,s)$ in the normalized critical strip
\[
            \frac w2<\Re s<\frac w2+1
\]
lies on $\Re s=(w+1)/2$. The same statement holds for the completed function.
For $M(n)$ the line is $\Re s=(w+1)/2-n$.
The assertion is a statement about divisors and does not assert
holomorphy on the critical line.

\paragraph{Proof.}
Anderson, Corollary~5.7.2 and Remark~5.7.3, printed p.~181
(PDF p.~30), supply number fields $F_a$, algebraic Hecke characters
$\psi_a$ of the \emph{same weight $w$}, and integers $m_a$, such that
\begin{equation}\label{eq:cm-factorization}
 L(M,s)=\prod_a L_{F_a}(\psi_a,s)^{m_a}.
\end{equation}
This is the consequence of Brauer induction for the Taniyama group,
not an assumption that the CM endomorphisms descend to $\mathbf Q$.
The equality uses the canonical local factors, including ramified places;
strict compatibility is Corollary~5.7.1. The same induced-representation
identity supplies the archimedean factors (equivalently,
\cite[\S7.1, Theorem~7.1; \S7.2, Corollaries~7.6 and~7.9]{FarguesCM}).
With a fixed coefficient-field embedding into $\mathbf C$, the corresponding
signed degree identity is
$\operatorname{rank}M=\sum_a m_a[F_a:\mathbf Q]$; if the coefficient
field is retained rather than restricted to $\mathbf Q$, this is understood
for the chosen coefficient component.

Use the ideal-character convention of Anderson, \S2.1, printed p.~159,
so that
$|\psi_a(\mathfrak p)|=(N\mathfrak p)^{w/2}$ at unramified primes.
Define the unitary normalization by
\[
 \chi_a(\mathfrak p)=\psi_a(\mathfrak p)(N\mathfrak p)^{-w/2}.
\]
Equivalently this is the corresponding norm twist of the idelic character.
For an algebraic Hecke character of weight $w$ the twist is a continuous
unitary character, including at the infinite places. One can see the
unitarity directly from its algebraic infinity type: the sum of a pair of
complex-conjugate exponents is $w$; at a real place the exponent is $w/2$.
After removing this real norm exponent its absolute value is trivial on
the archimedean connected group. The remaining absolute value factors
through the finite ray-class quotient and is therefore trivial there too.

Comparison of local Euler factors gives
\[
 L_{F_a}(\psi_a,s)=L_{F_a}(\chi_a,s-w/2).
\]
For $s_0$ in the stated strip off its center, put $z_0=s_0-w/2$.
Then $0<\Re z_0<1$ and $\Re z_0\ne1/2$.
By \textbf{(M)}, each completed Hecke factor is nonzero there. Classical
Hecke continuation gives no poles in this open strip; the possible poles
for a norm character are on its boundary. The unitary archimedean gamma
factors are finite and nonzero when $0<\Re z<1$, since their arguments
have positive real part. Consequently the finite Hecke factors are also
holomorphic and nonzero at $z_0$. Equation~\eqref{eq:cm-factorization}, with
its integer exponents, is therefore holomorphic and nonzero at $s_0$.
The identity initially holds in an absolute-convergence half-plane and
extends meromorphically, so this local argument proves the assertion for
the global function. In divisor notation it says
\[
 \operatorname{ord}_{s_0} L(M,s)
   =\sum_a m_a\operatorname{ord}_{z_0}L_{F_a}(\chi_a,z)=0.
\]
No cancellation can create a divisor away from the supports of the
individual factors. Applying $L(M(n),s)=L(M,s+n)$ proves the twist
assertion. \hfill$\square$

\paragraph{Scope and limits of CM01.}
The theorem applies to CM elliptic curves, all cohomology of CM abelian
varieties and their products, their tensor constructions and summands,
Artin twists of these objects, and the resulting descents to $\mathbf Q$.
For an existing Chow motive whose absolute-Hodge realization lies in this
category, its realized $L$-function has the same conclusion. It does not
assert that every absolute Hodge projector is algebraic.
For induced rank-one Hecke motives the factorization is an actual product
and holomorphy in the strip is already classical. For arbitrary CM motives
negative $m_a$ are essential: Artin motives themselves belong to
$\mathsf{PCM}$, so a claim of general holomorphy here would include
Artin's conjecture. Meromorphic divisor RH closes a different gate.
The complete function may have the usual endpoint divisors, and the finite
function has gamma-forced trivial divisors outside the displayed strip;
neither is being called a nontrivial zero in CM01.

\paragraph{The remaining gate is not determined by the weight alone.}
The Hecke reduction works because the connected Taniyama group is a
protorus. It does not express a motive with genuinely nonabelian connected
geometric monodromy as a virtual sum of induced rank-one characters.
The first effective examples are $h^1(E)$ for a non-CM elliptic curve and
the motives of non-CM modular eigenforms of weight at least two.
Tate twisting preserves this obstruction: for example,
$\operatorname{Sym}^2h^1(E)(1)$ has weight zero but remains non-CM.
Thus ``all remaining motives have positive weight'' is not an invariant
formulation. What remains is analytic continuation where unavailable,
and the appropriate RH for the higher-rank non-CM geometric factors;
known modularity does not itself supply this RH.



% Prepared by /root/intake_audit, GPT-6 Astra Pro, 21 November 2026 (Europe/Berlin).
% A = theorem relative to the stipulated Alpoge input; U = already unconditional.
\section{Audit of conditional statements and explicit arithmetic consequences}
\label{sec:intake-audit}

The audit distinguishes a source statement from an analytic estimate that one
may append to it.  In particular a previously proved existence theorem does
not become newly unconditional merely because an estimate of an associated
$L$-value improves.  In the tables, A means that the analytic hypothesis is
discharged by the stipulated Alp\"oge Dirichlet GRH, M means that the
additional general Hecke GRH of Masullo is used, U means previously
unconditional, and B means that a further hypothesis is still required.
``U0'' in the metadata inventory means only that no analytic RH hypothesis
occurs in the supplied row; it is not a new proof of that row.

\subsection{IA01: what the supplied conditional inventory actually contains}

\paragraph{Finding IA01 (checked document audit).}
The supplied \path{BRIDGE_FLAGS.tsv} contains 575 data rows: 456 under
motives/Tate twists, 64 under purity/weights, 53 under Weil conjectures, and
only two under RH/GRH.  The latter are the following.
\begin{enumerate}
\item S6\_SERIES\_GRIND, \path{maps:src:F.2.36}, source line 280:
the available context explicitly says that the class and unit groups were
certified with \path{GRH=false}.  This is not an unmet GRH hypothesis.
\item LM8\_PUSH, A.69, source line 533: the available context contains
``DS/RH'' in a discussion of genus, finite-flat fixed points and a tree.
Deuring--Shafarevich/Riemann--Hurwitz is the natural expansion in that
context; the source proof itself is absent, so this expansion is an
inference, not a checked quotation.  No analytic zero hypothesis can be
extracted from the supplied row.  It is not upgraded using either input.
\end{enumerate}
The four intakes contain no occurrence of ``Siegel''.  The SA07 intake
contains no Faltings-height or Colmez statement.  Those height questions
are therefore treated below as questions explicitly added by the prompt,
not assigned invented SA07 statement numbers.  A reproducible audit of
all 575 flag rows, all 587 identified Markdown table records, and the 51
analytic-keyword matches accompanies the paper.

\paragraph{Check.}
The audit parsed the TSV with its actual header, retained the original
statement and context in every row, counted categories, and searched the
four complete intake texts for GRH, Siegel, Littlewood, least non-residue,
and Riemann-hypothesis phrases.  The mathematical interpretation of every
positive hit was then read in its surrounding paragraph.  This is an
exhaustive audit of the supplied metadata and intake prose, not of the
unattached original research returns.

\subsection{IA02: explicit Dirichlet values, class numbers, and prime witnesses}

Put, for $q\geq 10^{10}$,
\begin{align*}
 U(q)&=2e^\gamma\left(\log\log q-\log2+\frac12+
              \frac1{\log\log q}\right),\\
 V(q)&=\frac{12e^\gamma}{\pi^2}
 \left(\log\log q-\log2+\frac12+
              \frac1{\log\log q}+\frac{14\log\log q}{\log q}\right).
\end{align*}

\paragraph{Finding IA02 (A).}
For every nonprincipal primitive Dirichlet character of conductor
$q\geq10^{10}$,
\[
             V(q)^{-1}\leq |L(1,\chi)|\leq U(q).
\]
For every negative fundamental discriminant $-D$ with $D\geq10^{10}$,
\[
 \frac{\pi\sqrt D}{12e^\gamma
 (\log\log D-\log2+\frac12+(\log\log D)^{-1}
                     +14\log\log D/\log D)}
 \leq h(-D)\leq\frac{\sqrt D}{\pi}U(D).
\]
For every prime $p\geq5$, the least positive quadratic non-residue
$n(p)$ satisfies $n(p)<(\log p)^2$.  For every $q>3$ and
$(a,q)=1$, the least prime $P(a,q)$ satisfies
$P(a,q)\leq(\varphi(q)\log q)^2$.

\paragraph{Proof.}
Lamzouri--Li--Soundararajan, \emph{Conditional bounds for the least
quadratic non-residue and related problems}, corrected arXiv version
1309.3595v3, Theorem 1.5 (pp.~3--4), Corollary 1.3 (p.~4),
Corollary 1.1 (p.~2), and Corollary 1.2 (p.~3), have precisely these
Dirichlet GRH hypotheses.  The stipulated input supplies them.
The class-number constants also follow directly by multiplication by
$\sqrt D/\pi$ in the quadratic class-number formula (their equation
(1.3)).  The non-residue is prime: if $n(p)$ were composite, every
prime divisor would be smaller than $n(p)$ and hence a residue by
minimality; their product with multiplicities would also be a residue,
a contradiction.  No extension of the stated conductor threshold
is asserted.  The corrected version is essential: its asymptotic
coefficient for the quadratic non-residue is $0.8$, not the $0.42$
appearing in an obsolete draft; that asymptotic coefficient is not used
here. \hfill$\square$

For abelian $K/\mathbb Q$ one obtains a completely explicit version for
$h_KR_K$.  Let $X_K$ denote the primitive Dirichlet characters attached
to its characters, let $r_1,r_2,w_K,D_K$ have their usual class-number
formula meanings, and choose rigorously certified positive lower and
upper bounds $\ell_\chi,u_\chi$ for $|L(1,\chi)|$.  For conductors at
least $10^{10}$ one may use $\ell_\chi=V(q_\chi)^{-1}$ and
$u_\chi=U(q_\chi)$; smaller conductors must be evaluated or bounded
separately.  Then
\[
 \frac{w_K\sqrt{|D_K|}}{2^{r_1}(2\pi)^{r_2}}
 \prod_{\chi\in X_K,\ \chi\ne1}\ell_\chi
 \leq h_KR_K\leq
 \frac{w_K\sqrt{|D_K|}}{2^{r_1}(2\pi)^{r_2}}
 \prod_{\chi\in X_K,\ \chi\ne1}u_\chi.
\]
This is the abelian zeta factorization followed by the analytic
class-number formula (A.~V.~Sutherland, \emph{18.785 Number Theory I},
Fall 2021, Lecture 19, Theorems 19.15 and 19.12, pp.~11 and 9);
conjugate factors pair to positive real numbers.
A bound for $h_KR_K$ is not a bound for $h_K$ until a positive lower
bound for $R_K$ has been supplied.  In particular the SA06 B.4e
exponent involving $h_K$ needs that additional step.  Exact algorithms
for a fixed field or an independent regulator bound provide it; RH
alone does not turn a product into either of its factors.

\subsection{IA03: which supersingular count receives the square-root bound}

\paragraph{Finding IA03 (A, with an unconditional counting identity).}
Let $S_p$ be the number of supersingular $j$-invariants in
$\mathbb F_p$, where $p>3$.  Put
\[
 (D_p,c_p)=
 \begin{cases}
 (4p,\tfrac12),&p\equiv1\pmod4,\\
 (p,1),&p\equiv7\pmod8,\\
 (p,2),&p\equiv3\pmod8.
 \end{cases}
\]
If $D_p\geq10^{10}$, then
\[
 \frac{c_p\sqrt{D_p}}{\pi V(D_p)}
 \leq S_p\leq\frac{c_p\sqrt{D_p}}{\pi}U(D_p).
\]
\paragraph{Proof.}
The identity $S_p=c_ph(-D_p)$ is Delfs--Galbraith,
\emph{Computing isogenies between supersingular elliptic curves over
$\mathbb F_p$}, arXiv:1310.7789, equation (1), p.~2.
In each case $-D_p$ is a fundamental discriminant.  Substitution in
IA02 proves the two inequalities. \hfill$\square$

This resolves the proposed estimate for SA06 A.8a.  It does not estimate
the total number of geometric supersingular $j$-invariants: that number
is $\lfloor p/12\rfloor+\epsilon_p$ with
$\epsilon_p=0,1,1,2$ in the four residue classes modulo 12
(Delfs--Galbraith, p.~1, citing Silverman V.4.1(c)).
SA06 A.8c only supplies the phrase ``maximal-order trace-zero triples
have class number $h(R)$''; its definition of $R$ and full count are not
attached.  The neighboring A.1f row explicitly mentions a maximal
quadratic order, making that the likely intended setting.  The precise
transfer criterion is: if $R$ is the maximal order of an imaginary
quadratic field of discriminant $-D$, then IA02 applies to $h(R)$ with
that $D$; for $R=\mathcal O_{\mathbb Q(\sqrt{-p})}$, use $D=p$
when $p\equiv3\pmod4$ and $D=4p$ otherwise.  Identifying this $R$
from the original A.8c proof is the remaining source gate.  A quaternion
ideal-class count, if that had instead been intended, would require its
different mass formula.
The archive locator carried by the intake for the A.8 question is cote
15, scan 70; that page is reported to end mid-question.  No scan was
provided in this package, and no new transcription reading is claimed.

\subsection{IA04: effective CM elliptic heights with a fixed normalization}

Use the Deligne normalization of the stable Faltings height in
Habegger--Jones--Masser, \emph{Six unlikely intersection problems in
search of effectivity}, equation (5.9) and Lemma 14(i) (printed p.~22).
Thus, for CM by the order of discriminant $-Df^2$, with $-D$ fundamental,
\[
 h_F(E)=\frac14\log(Df^2)+\frac12\frac{L'}L(1,\chi_{-D})
 -\frac12\sum_{p\mid f} e_f(p)\log p
 -\frac12(\gamma+\log(2\pi)),
\]
where, for $p^a\Vert f$,
\[
 e_f(p)=\frac{1-\chi_{-D}(p)}{p-\chi_{-D}(p)}
              \frac{1-p^{-a}}{1-p^{-1}}.
\]
This identity is unconditional and includes the nonmaximal-order
correction; GRH is used only for the next estimate.

\paragraph{Finding IA04 (A).}
Let $D>e^6$, set $t=\tfrac12\log D$ and
$b=\gamma+\log(2\pi)-2$, and define
\[
 R_-(D)=\frac{t-1}{t-3}\left[
 2(\log t+1)-\frac{b-\log2}{t-1}
             +\frac{\log2+1/4}{t^2-1}\right].
\]
Then
\[
 \left|h_F(E)-\frac14\log(Df^2)
       +\frac12\sum_{p\mid f}e_f(p)\log p
       +\frac12(\gamma+\log(2\pi))\right|
                    <\frac12R_-(D).
\]
In particular, for maximal CM orders,
$h_F(E)=\tfrac14\log D+O(\log\log D)$ with the displayed effective
error, once the fixed additive normalization is included.

\paragraph{Proof.}
The character $\chi_{-D}$ is primitive, odd, nonprincipal and has
conductor $D$.  Ihara--Murty--Shimura, \emph{On the logarithmic
derivatives of Dirichlet $L$-functions at $s=1$}, Corollary 3.3.2,
equation (3.3.3), p.~9, gives
$|L'/L(1,\chi_{-D})|<R_-(D)$ under GRH.  Here $t>3$, so multiplying
by $(t-1)/(t-3)$ preserves the inequality.  Substitute this in the
unconditional height identity.  This also verifies the factor $1/2$
and every additive constant.  The range $D\leq e^6$ is not silently
absorbed into the stated bound; its exact height formula remains
available. \hfill$\square$

\subsection{IA05: the no-Siegel-zero gate is smaller than general Hecke GRH}

\paragraph{Finding IA05 (A plus Colmez where still needed).}
Let $(E,\Phi)$ be a CM type of degree $[E:\mathbb Q]=2g$ and let
$A/K$ have multiplication $\mathcal O_E\hookrightarrow
\operatorname{End}_K(A)$ of this type.  Assuming the individual
Colmez formula for $(E,\Phi)$, the stipulated Dirichlet GRH gives
\[
 h_F^{\rm st}(A)<
 \frac{75}{4}g((2g)!)^2\log\bigl(|D_K|^{1/[K:\mathbb Q]}\bigr)
 -\frac g4(\gamma+\log(4\pi)).
\]
For abelian CM fields the Colmez formula is already a theorem, so this
is an A consequence.  For general CM types the individual Colmez
formula remains a separate special-value gate; Masullo's zero claim
does not prove it.

\paragraph{Proof.}
Dirichlet GRH excludes zeros in the rectangle
$1-(4\log d)^{-1}\leq\Re s<1$,
$|\Im s|\leq(4\log d)^{-1}$ for every imaginary quadratic
discriminant of absolute value $d\geq3$; its left edge is greater
than $1/2$.  Thus the constant $C_{\rm zero}=4$ is available in
Wei's Conjecture 1.1.  For the CM Galois closure of the reflex field,
Stark's descent of an exceptional zero to an imaginary quadratic
subfield, as applied in Wei, \S4.3, proof of Theorem 1.3 (p.~18),
excludes such a zero in the full rectangle with constant $1/4$.
Wei Proposition 4.7 (p.~17) gives the required zero exclusion for
the odd Artin characters in the Colmez expression.
Apply Wei Corollary 4.3, equation (13), p.~14, with $c=1/4$ and
therefore $c'=0$.  It yields
\[
 h_F^{\rm st}(A)<\frac{75}4g(2g)!\log|D_{E^*}|
                         -\frac g4(\gamma+\log(4\pi)).
\]
Because all specified CM endomorphisms are defined over $K$, the
reflex field $E^*$ embeds in $K$.  Discriminant transitivity gives
$\log|D_{E^*}|/[E^*:\mathbb Q]\leq
\log|D_K|/[K:\mathbb Q]$, and $[E^*:\mathbb Q]\leq(2g)!$.
Substitution proves the claim.  The abelian case uses Wei Theorem
4.8 (p.~18), recording Colmez--Obus. \hfill$\square$

The numerical coefficient in the general bound is deliberately coarse
but explicit; for $g=1$ it is 75 and for $g=2$ it is 21600.
This is an extraction from the proof, not an invented value for an
unspecified effective constant.  It is stronger as a logical reduction
than saying that every CM no-Siegel-zero hypothesis must await Masullo.
Wei Proposition 5.3 additionally combines the averaged Colmez theorem
with a universal lower height bound to obtain discriminant-of-$E$
bounds without the individual Colmez formula.  Its universal lower
height constant has not been numerically re-estimated here.

\subsection{IA06: the cyclotomic quartic height is already unconditional}

\paragraph{Finding IA06 (U; an explicit form of the known Colmez formula).}
For a CM abelian surface with full multiplication by
$\mathbb Z[\zeta_5]$, the same height normalization gives
\[
 h_F^{\rm st}(A)=\frac12\log5
 -\frac32\log\frac{\Gamma(1/5)}{\Gamma(4/5)}
 -\frac12\log\frac{\Gamma(2/5)}{\Gamma(3/5)}
 =-1.4525092396456446503\ldots.
\]
Neither zero theorem is needed for this identity.

\paragraph{Proof.}
Let $G=C_4=\langle\sigma\rangle$ and choose the CM type
$\{1,\sigma\}$; all four CM types are conjugate.  The class function
$A^0_\Phi$ of Wei's definition (p.~2) has values $(2,1,0,1)$.
If $\chi(\sigma)=i$, its character decomposition is
$1+\tfrac12\chi+\tfrac12\bar\chi$.
The unconditional abelian Colmez formula (Wei Theorem 4.8)
therefore gives
$h_F=-\Re(L'/L)(0,\chi)-\tfrac12\log5$:
the trivial-character term cancels the prescribed $-\log(2\pi)$
term in $Z_\Phi$.
Choose $\sigma$ corresponding to 2 modulo 5, so
$\chi(1),\chi(2),\chi(3),\chi(4)=(1,i,-i,-1)$.
The finite Hurwitz-zeta expansion and its values at zero
(NIST DLMF, equations 25.15.3, 25.11.13 and 25.11.18) yield
\[
 L(0,\chi)=\frac{3+i}{5},\qquad
 \frac{L'}L(0,\chi)=-\log5+
 \frac{5}{3+i}\left(
 \log\frac{\Gamma(1/5)}{\Gamma(4/5)}+
 i\log\frac{\Gamma(2/5)}{\Gamma(3/5)}\right).
\]
Taking real parts gives the formula.  The decimal is a numerical
evaluation, not a replacement for the exact expression.
Moreover the unique quadratic subfield of $\mathbb Q(\zeta_5)$ is
$\mathbb Q(\sqrt5)$, which is real; hence Wei Theorem 1.4 already
gave its qualitative root-discriminant height bound before either
stipulated input. \hfill$\square$

\subsection{IA07: the $23$-cyclotomic obstruction and its exact certificate}

\paragraph{Finding IA07 (U).}
The arithmetic class group used in the SA20 A5.011 bridge has
$h(\mathbb Q(\zeta_{23}))=3$, and Rim's theorem identifies
$\widetilde K_0(\mathbb Z[C_{23}])\simeq\mathbb Z/3\mathbb Z$.
Its relative class-number computation
is unconditional and uses $L(0,\chi)$, not GRH.  Bernoulli numbers
alone compute $h^-$; the plus part must also be supplied.

\paragraph{Proof.}
Use 5 as a primitive root modulo 23 and let
$S(x)=\sum_{j=0}^{21}a_jx^j$, where $a_j$ is the least positive
residue of $5^j$ modulo 23.  The odd characters have $\chi(5)$ equal
to the eleven roots of $x^{11}+1$.  Reduce $S$ modulo that polynomial:
\[
\begin{split}
 T(x)={}&-21-13x-19x^2-3x^3-15x^4+17x^5-7x^6\\
        &+11x^7+9x^8-x^9-5x^{10}.
\end{split}
\]
The determinant of multiplication by $T$ in
$\mathbb Z[x]/(x^{11}+1)$ is
\[
                    -127262242448329728=-3072\cdot23^{10}.
\]
The attached exact integer-arithmetic script gives the matrix and
computes its determinant by fraction-free elimination, independently
checking it against a symbolic determinant.  Since
$L(0,\chi)=-S(\chi(5))/23$, this yields
$\prod_{\chi\ {\mathrm{odd}}}L(0,\chi)=3072/23$.
The relative class-number formula
$h^-=2p\prod_{\chi\ {\mathrm{odd}}}(-B_{1,\chi}/2)$
(Washington, \emph{Introduction to Cyclotomic Fields}, Theorem 4.17;
also displayed in D.~S.~Thakur, \emph{Fermat's last theorem for regular
primes}, in the \emph{Cyclotomic Fields} lecture collection, p.~172)
now gives $h^-=46\cdot3072/(23\cdot2^{11})=3$.
Miller, \emph{Real cyclotomic fields of prime conductor and their
class numbers}, arXiv:1407.2373, Theorem 1.1, gives $h^+=1$ for
prime conductor at most 151, unconditionally.  Thus $h=h^+h^-=3$,
and a group of order 3 is cyclic. \hfill$\square$

For the final identification one uses the known theorem of
Rim, \emph{Modules over finite groups}, Theorem 6.24, p.~711.
Its exact statement and locator are confirmed in Ullom,
\emph{Normal bases in Galois extensions of number fields}, p.~159;
the isomorphism of locally free class groups is also Siviero,
\emph{Realisable classes, Stickelberger subgroup and its behaviour
under change of the base field}, Lemma 4.5, equation (16), p.~82.
This verifies the arithmetic value and the stated $K$-theoretic bridge.
The topological application is separate; its original A5.011 proof
is not attached.  In particular the intake's
shorthand $L(1,\chi)\asymp\log\log q$ must be replaced by
$V(q)^{-1}\leq|L(1,\chi)|\leq U(q)$, not by a two-sided estimate
with $\log\log q$ on both sides.

\subsection{IA08: the $5$-cyclotomic unit and what its bound does not prove}

\paragraph{Finding IA08 (U for the value, A for family estimates).}
For $\phi=(1+\sqrt5)/2$,
\[
             L(1,\chi_5)=\frac{2\log\phi}{\sqrt5},\qquad
             \phi^4=\frac{7+3\sqrt5}{2}.
\]
Thus the numerical absolute value reported by REVLOC I.8b is
$6.8541019662496845\ldots$.  This is not a proof that an unspecified
group-ring unit is a generator of $\operatorname{Wh}(C_5)$; the
intake says that the application only needs a nontrivial unit.
Its status table nevertheless marks I.5 ``refuted as worded'';
this new arithmetic check does not reverse that original verdict.

\paragraph{Proof.}
The elementary Fourier formula for a primitive even character gives
\[
 L(1,\chi_5)=-\frac1{\sqrt5}
       \sum_{a=1}^4\chi_5(a)\log(2\sin(\pi a/5)).
\]
For completeness this form follows by expanding
$-\log|1-re^{2\pi ia/5}|=\sum_{n\geq1}r^n
\cos(2\pi na/5)/n$ for $0<r<1$, summing with the weights
$\chi_5(a)$, using the quadratic Gauss sum $\sqrt5$, and taking
the Abel limit $r\uparrow1$; the Dirichlet series converges at 1.
The two equal sine pairs turn the result into
$(2/\sqrt5)\log(\sin(2\pi/5)/\sin(\pi/5))$.
The sine ratio is $2\cos(\pi/5)=\phi$.
The formula for $\phi^4$ follows from $\phi^2=\phi+1$.
The identities are unconditional.  The Littlewood bounds of IA02
apply to families of real quadratic class-number--regulator products;
they neither create this particular unit nor prove an integral
Whitehead-group generator statement. \hfill$\square$

\subsection{Precise effect on the intake rows}

\paragraph{Finding IA09 (a counterexample family to a uniform
$\log\log q$ lower bound).}
The SA20 phrase $L(1,\chi)\asymp\log\log q$ cannot be repaired just
by adding absolute-value signs.  For every odd prime $p$ there is a
nonprincipal primitive odd character $\chi_p$ modulo $p$ such that
$|L(1,\chi_p)|\leq\pi/\sqrt6$.  This supplies an infinite family
contradicting any positive uniform lower bound
$c\log\log p$ for all primitive characters.

\paragraph{Proof.}
For an odd character modulo $p$, the elementary finite Fourier
formula at 1 gives
\[
 |L(1,\chi)|^2=\frac{\pi^2}{p^3}
              \left|\sum_{a=1}^{p-1}a\bar\chi(a)\right|^2.
\]
To check this particular form, multiply the Abel limit of
$\sum_{n\geq1}\sin(2\pi na/p)/n=\pi(1/2-a/p)$ by
$\bar\chi(a)$ and sum over $a$; the Gauss sum has absolute value
$\sqrt p$ and the constant term vanishes.  Orthogonality over the
$(p-1)/2$ odd characters says
\[
 \sum_{\chi\ {\mathrm{odd}}}\bar\chi(a)\chi(b)
       =\frac{p-1}{2}\bigl(\mathbf1_{a=b}-\mathbf1_{a=p-b}\bigr).
\]
Consequently their mean squared $L$-value equals
\[
 \frac{\pi^2}{p^3}
 \left(\sum_{a=1}^{p-1}a^2-\sum_{a=1}^{p-1}a(p-a)\right)
       =\frac{\pi^2(p-1)(p-2)}{6p^2}<\frac{\pi^2}{6}.
\]
Take a character with value no greater than the mean.  It is
nonprincipal because it is odd, and every nonprincipal character
modulo a prime is primitive.  Infinitely many primes exist and
$\log\log p$ is unbounded.  No RH assertion is used.\hfill$\square$

SA06 A.8a receives IA03; its underlying count remains U.  For A.8c
the missing definition of $R$ prevents a narrower estimate.  SA06
B.4a--B.4e retain their respective original algebraic statuses; an
abelian-field $hR$ estimate is A and a full nonabelian Dedekind-GRH
estimate is M, while a bound on $h$ requires regulator control.
SA07 E.2.1--E.2.6 and the other integral-Betti rows do not assume
zero statements: the new RH statement for their nonabelian ring-class
field zeta functions is M, with no change to the place-separation
proofs.  SA07 D.3.2's open-image theorem is U; its proposed effective
nonabelian division-field Chebotarev input is M.  The A.4 period
values remain special-value and transcendence questions; locating
the associated Hecke zeros is M but is not a proof of their period
claims.  SA20 A5.011 and REVLOC I.5/I.8b were already U in their
stated arithmetic uses; REVLOC I.5 itself remains refuted as worded,
as its intake explicitly records.  None of these conclusions supplies a new
archive transcription: the supplied files contain second-hand
locators, not the original scans.


\section{Special values, extensions, and what zero location does not compute}
\label{sec:special-values}

\subsection{Version-specific Deligne--Goncharov citations}
All citations marked DG-att below refer to the supplied February 2003 source,
not silently to a later pagination.  DG-att \S1.6 is in
\path{Sec1ccc.tex}; its Proposition~1.7 characterizes unramified mixed Tate
motives by their \(\ell\)-adic realizations, and Proposition~1.9 states
invariance of the higher \(\operatorname{Ext}^1\)'s and vanishing of Yoneda
\(\operatorname{Ext}^2\).  In \path{Sec2ccc.tex}, equation (2.1.3) gives
\[
 \operatorname{Ext}^1_{\mathrm{MT}(\mathcal O_S)}
   (\mathbb Q(0),\mathbb Q(n))=
 \begin{cases}
 0,&n\leq0,\\
 \mathcal O_S^\times\otimes\mathbb Q,&n=1,\\
 K_{2n-1}(k)\otimes\mathbb Q,&n\geq2.
 \end{cases}
\]
Propositions~2.2--2.3 describe the graded prounipotent Lie algebra and its
freeness.  The warning in \S2.4 says the choice of free generators is not
canonical.  DG-att \S2.17 and Proposition~2.18 describe the descent category
\(\mathrm{MAT}(k)\) and its mixed Tate subcategory.  The 91-page IAS-hosted
\path{Tate.pdf} inspected for comparison has additional material: its
\S1.7, not its \S1.6, defines \(\mathrm{MT}(\mathcal O_S)\).  This is version
drift, not evidence of a transcription error in Grothendieck's manuscripts.

\paragraph{Existing arithmetic content.}
Borel, \emph{Stable real cohomology of arithmetic groups}, Ann. Sci. ENS
\textbf{7} (1974), Proposition~12.2, p.~271, gives for a number ring with
\(r_1\) real and \(r_2\) complex places, in degrees \(j\geq2\), the ranks
\(0,r_1+r_2,0,r_2\) according as \(j\equiv0,1,2,3\pmod4\).
Thus for \(\mathrm{MT}(\mathbb Z)\), (2.1.3) already gives a
one-dimensional extension space in each odd degree \(n\geq3\), and zero in
the other positive degrees.  Both this rational statement and DG-att's free
Lie-algebra conclusion precede and do not require \(A\) or \(M\).
They do not calculate the torsion of integral \(K_j(\mathbb Z)\).
Tensoring with \(\mathbb Q\) sends every finite abelian group to zero, so a
rational vanishing statement cannot distinguish \(0\) from, for example,
\(\mathbb Z/3\mathbb Z\).

\subsection{Finding SV01: the Euler factor forgets extension classes}
\label{SV01}
\begin{proposition}[SV01]
Let \(V\) be a finite-dimensional \(\ell\)-adic representation, unramified
at a prime \(p\), with a finite Galois-stable filtration.  Then
\[
 \det(1-\operatorname{Frob}_pT\mid V)
 =\prod_j\det(1-\operatorname{Frob}_pT\mid\operatorname{gr}_jV).
\]
Consequently, whenever complex local factors and an Euler product absolutely converging
in a right half-plane are defined, the incomplete product \(L^S(V,s)\),
with all ramified places in the finite set \(S\), and every zero-location
statement about it depend only on
the class \([V]=\sum_j[\operatorname{gr}_jV]\) in the Grothendieck group.
In particular, the mixed Tate/Artin--Tate transfer theorems do not determine
extension classes or the period matrix of an object.
\end{proposition}
\begin{proof}
Choose a basis adapted to the finite filtration.  The Frobenius matrix is
block triangular with diagonal blocks the induced Frobenius matrices on the
graded pieces.  The determinant of \(I-TF\), where \(F\) is this matrix, is the product
of the corresponding diagonal-block determinants.  Multiplying their
inverses at \(T=p^{-s}\) gives the Euler-product identity wherever the
products converge absolutely.  Any meromorphic continuation of one side is
therefore the continuation of the other side.  Extension parameters occur
only in off-diagonal blocks and cannot affect this polynomial.  This proves
the asserted dependence and the limitation.  We deliberately exclude
ramified places: passage to inertia invariants need not preserve an exact
sequence with unipotent inertia, so this argument by itself does not establish
multiplicativity of every possible complete ramified Euler-factor convention.
\end{proof}

This is the determinant realization of the identity
\(\lambda_t([V])=\prod_j\lambda_t([\operatorname{gr}_jV])\).
Its Grothendieck formalism is SGA~6, Expos\'e~V, \S2 (definition and examples
of pre-\(\lambda\)-rings), and \S7, equation (7.1.1) (Adams operations).
Equivalently, for a fixed Frobenius endomorphism \(F\),
\[
 -\log\det(1-FT)=\sum_{r\geq1}\operatorname{Tr}(F^r)T^r/r.
\]
The trace package therefore retains Frobenius eigenvalues and weights but
forgets extension parameters.  No EGA or FGA result is being claimed to prove
an analytic zero theorem.

\subsection{Finding SV02: an actual mixed Tate counterexample to recovery}
\label{SV02}
\begin{proposition}[SV02]
There are nonisomorphic mixed Tate motives over \(\mathbb Z[1/2]\) with the
same incomplete \(L\)-function, hence the same nontrivial zero set.  Their
extension classes are different, before either RH input is used.
\end{proposition}
\begin{proof}
By DG-att (2.1.3),
\(\operatorname{Ext}^1(\mathbb Q(0),\mathbb Q(1))=
\mathbb Z[1/2]^\times\otimes\mathbb Q\).
The valuation at 2 sends the class of the unit 2 to 1; it is therefore
nonzero.  Let \(E_2\) be the corresponding Kummer extension, described in
the proof of DG-att Proposition~1.7, and let
\(E_0=\mathbb Q(0)\oplus\mathbb Q(1)\).
The weight filtration is intrinsic and functorial, by DG-att \S1.1.
An isomorphism between these objects would induce invertible scalars on the
two one-dimensional graded pieces and would carry the extension class of
\(E_2\) to zero.  Multiplication by nonzero scalars cannot annihilate a
nonzero vector in the rational extension space.  Thus the objects are
nonisomorphic.  For each \(p\neq2,\ell\) the realizations are unramified,
and geometric Frobenius has eigenvalues \(1,p^{-1}\) on their graded pieces.
For \(S\) containing \(2\) and \(\ell\), SV01 gives
\[
 L^S(E_2,s)=L^S(E_0,s)=\zeta^S(s)\zeta^S(s+1).
\]
The identity is exact and precedes RH.  It proves non-recovery of the
extension class from this \(L\)-function, not a formal independence theorem
between arbitrary conjectures.
\end{proof}

\subsection{Bernoulli numbers and regulators}
The classical identity at negative integers is separate from a statement
about zeros in a critical strip.  In an unambiguous convention, put
\[
 B_{n,\chi}=f^{n-1}\sum_{a=1}^{f}\chi(a)B_n(a/f).
\]
Then for \(n\geq1\),
\(L(1-n,\chi)=-B_{n,\chi}/n\).  The exact references are the Hurwitz
expansion of a Dirichlet \(L\)-function, NIST DLMF (25.15.3), and
\(\zeta(-m,a)=-B_{m+1}(a)/(m+1)\), DLMF (25.11.14).  With \(f=1\), this
convention uses \(B_1(1)=1/2\), so also correctly gives \(\zeta(0)=-1/2\).
For \(m\geq1\), \(\zeta(1-2m)=-B_{2m}/(2m)\).
For the higher-regulator comparison one may cite Borel, \emph{Cohomologie
de SLn et valeurs de fonctions z\^eta aux points entiers}, Ann. Scuola
Norm. Sup. Pisa \textbf{4} (1977), Theorem~6.2 and \S6.4, pp.~631--633,
with its 1980 erratum, p.~373; the normalization of the regulator must be
retained.  These identities, analytic class-number formulas, and regulator
comparison theorems are additional arithmetic input, not consequences of moving
nontrivial zeros to a line.  GRH does improve effective estimates for some
special values; the bounds in the audit use those separate theorems.  It
does not replace exact Bernoulli arithmetic or determine a regulator matrix.

In particular, the SA20 intake's A5.011 cyclotomic-class obstruction and
A.103 low-degree \(K\)-groups are not newly proved by \(A\).
The finite \(\mathrm{Cl}(\mathbb Q(\zeta_{23}))\) computation cited there
and the uses of \(\zeta(-1)\) are logically distinct from an effective
family estimate for \(L(1,\chi)\).  The separate intake audit records their
precise source status.  Nothing here asserts that zero location decides
Vandiver, integral \(K\)-group vanishing, a canonical splitting of the
motivic Galois Lie algebra, or the period conjectures.

\paragraph{Crystals and divided powers.}
The note's step 2 is stated for ordinary iterated invariant derivatives.
If one passes to Taylor coefficients normalized by \(a!\), the valuation
changes by \(v_p(a!)=\sum_{r\geq1}\lfloor a/p^r\rfloor\).  A bound with
leading term \(a/p+a/p^3+\cdots=ap/(p^2-1)\) therefore supplies only the
lower bound \(-a/(p^2-1)-O(\log a)\) for that quotient; it does not by
itself imply its integrality.  This is a comparison of bounds, not a claim
that the actual theta coefficients fail to be integral.  A crystal with
specified divided-power structure would require its own integral model and
compatibility theorem.  DG-att's introduction explicitly excludes the
crystalline realization from the realizations then constructed.  The
attached material consequently supplies no new general theorem about
crystals or divided powers as a consequence of RH.



\section{Checked observations on the analytic route}


\paragraph{Finding MDOC01 (the intended RH assertion is not a completed
numbered theorem in the supplied draft).}
The source convention for \textbf{(M)} is important. At PDF p.~522,
Theorem~13.3.1.2, headed ``Divisor of $L(s,\chi)$'', has no mathematical
statement or proof. Its footnotes request insertion of the statement and
proof of the main theorems. The preceding line refers to ``Theorem ??''.
At PDF p.~528, \S14.1.2 describes the intended consequence as critical-line
location and a noncentral multiplicity bound $N(\rho)\le[F:\mathbf Q]$,
while carrying a correction footnote. Accordingly CM01 is conditional on
the stated claim \textbf{(M)}, not a claim that this run certified the draft.
\emph{Verification:} direct reading of these pages, including a rendered
image of PDF p.~522, establishes this documentary finding.


\paragraph{Finding ML01 (checked local-pole separation, independent of RH).}
Let $F$ be a number field and $\chi$ a continuous unitary Hecke character.
The local factors at two finite places of distinct residue characteristic
have at most one common pole. If $\chi$ has finite order, then any fixed
$s\ne0$ is a pole of the local factor at at most $[F:\mathbf Q]$ finite
places. Under this finite-order hypothesis, $0$ is the only possible
point at which infinitely many finite local factors have poles. This independently verifies the content
of Masullo's Theorem~1.7.1.1(1),(3), PDF p.~35, in a slightly quantified form.

\paragraph{Proof.}
A ramified factor is $1$, so only unramified factors matter. Such a factor
is $(1-\chi_v(\pi_v)q_v^{-s})^{-1}$. If $s$ and $t$ are two distinct common
poles at $v_1,v_2$, division of the two pole equations gives
$(s-t)\log q_{v_j}\in2\pi i\mathbf Z$ for $j=1,2$.
Both integers are nonzero, so $\log q_{v_1}/\log q_{v_2}\in\mathbf Q$.
Writing $q_{v_j}=p_j^{f_j}$ contradicts unique factorization when
$p_1\ne p_2$. This proves the first assertion.
For the second, let $d$ be the order of $\chi$. At a pole,
$\exp(-ds\log q_v)=1$. If $s\ne0$ is a pole at two places of different
residue characteristic, this again makes the two logarithms rationally
dependent. Thus all the possible places lie over one rational prime, and
there are at most $[F:\mathbf Q]$ of them. \hfill$\square$

\paragraph{Finding ML02 (the restricted trace domain contains no nonzero
standard positivity test).}
Let $G$ be a locally compact abelian group with Haar measure and let
$H\subset G$ contain its identity. If $\psi$ is a continuous compactly
supported function, put $\psi^*(g)=\overline{\psi(g^{-1})}$.
If $(\psi*\psi^*)|_H=0$, then $\psi=0$.
Consequently the domain $\phi|_{C_F^1}=0$ in the draft's
Theorem~12.5.1.1 does not by itself furnish nonzero tests of the form used
in Weil's positivity criterion.
The draft itself states this limitation in \S1.7.2, PDF p.~35.
ML02 independently checks its general locally compact formulation; it is
not a new objection to a route that explicitly avoids that criterion.

\paragraph{Proof.}
At the identity,
$(\psi*\psi^*)(1)=\int_G|\psi(g)|^2\,d\mu(g)$.
Vanishing on $H$ makes this integral zero. Thus $\psi=0$ almost everywhere,
and continuity gives $\psi=0$ identically. \hfill$\square$

\paragraph{Check status.}
CM01, ML01 and ML02 have complete short proofs above. Their entire proofs
were read again in this run for normalization, hypotheses, and the
distinction between divisors and holomorphy. This is a human-readable
model proof check, not a formal certificate. The long functional-analytic
proofs behind the method map have not been certified in this run.


\section{Targeted verification of the Grothendieck archive}
\label{sec:archive-verification}

The supplied bridge paper gave text records for cote 34.  In this run the
Montpellier primary PDFs were accessed independently and the specific pages
below were rendered and read.  A \emph{PDF scan} counts the added cover as
page 1; the \emph{archival page} is the number printed beneath the facsimile.
These numbers must not be interchanged.  The source addresses are
\url{https://grothendieck.umontpellier.fr/34.pdf} and
\url{https://grothendieck.umontpellier.fr/60.pdf}.
Selected-page hashes and the inspection records accompany
the source table; the primary URLs identify the facsimiles.  No full-manuscript reading is claimed.

\subsection{The two Riemann--Weil passages in cote 34}

\paragraph{Scan 64, archival page 63, manuscript page II.12.}
The lower part of the page is headed \emph{Remarque 1.17}.  The remark is
crossed out by long diagonal strokes; it must not be silently promoted to
an uncancelled proposition.  The supplied transcription of its central
sentence agrees with the visible manuscript:
\begin{quote}
il est plausible que l'on ait encore
\(\overline{[\pi_0,H]}\) ouvert dans \(\pi'\).
Cela semblerait devoir r\'esulter de l'hyp. de Riemann--Weil en
dim. quelconque
\end{quote}
The continuation compares this with 1.1 and Riemann--Weil in dimension 1.
The words \emph{plausible} and \emph{semblerait} are substantive: the page
records a proposed deduction, not that deduction's proof.  The reference
is to geometric Riemann--Weil input.  Nothing in this inspection identifies
it with the critical-line theorem for \(\zeta(s)\) or a Dirichlet
\(L\)-function.  Q03 therefore retains its separate geometric-monodromy
gate; the inspection does not settle it.

\paragraph{Scan 74, archival page 73, manuscript page II.22.}
The label (2.5.5) is visible near the top.  Its following sentence agrees
with the supplied text record:
\begin{quote}
C'est une cons\'equence bien connue de l'hyp. de Riemann--Weil,
prouv\'ee en dim. 1.
\end{quote}
There is one narrow literal difference in the displayed formula.  The
manuscript writes the trivial invariant group as \(1\), whereas
\path{BRIDGES_20261120.tex} and \path{BRIDGES_QUESTIONS.tsv}
write it additively as \(0\).  The manuscript's preceding words say that
the last term is zero, and its subsequent reformulation on this same page
is
\[
 H^0\bigl(S,R^1f_*\mathbb Z_\ell\bigr)=0.
\]
Thus the literal correction is to record the original \(1\) and identify
the supplied \(0\) as additive normalization.  It does not change the
vanishing assertion or supply a new mathematical objection.  This is the
only transcription difference established by the targeted comparison.
The dimension-one Riemann--Weil attribution is already explicit in the
manuscript and is independent of assumption \(A\).

\subsection{Cote 60: determinant algebra and its geometric input}

The primary cover identifies cote 60 as \emph{Fonctions L}, with the
subtitle \emph{\'Equation fonctionnelle des fonctions L (cas g\'eom\'etrique)}
and 69 archival pages.  The PDF has the additional institutional cover,
so it has 70 scans.  Scan 2, archival page 1, is the manuscript folder
cover.  The following mathematical pages were read selectively:

\begin{center}
\begin{tabular}{>{\raggedright\arraybackslash}p{.12\textwidth}>{\raggedright\arraybackslash}p{.16\textwidth}>{\raggedright\arraybackslash}p{.61\textwidth}}
\toprule
PDF scan & Archival page & Visible content and usable scope \\
\midrule
3 & 2 & A finite-dimensional endomorphism, its characteristic polynomial,
exterior-power traces, and the reciprocal determinant expressed through
symmetric-power traces.  These are formal linear-algebra identities. \\
5 & 4 & Exterior and symmetric powers, and the corresponding
\(\lambda\)-operations on determinant series.  Several working formulas
are crossed out; no general arithmetic special-value theorem is asserted
here. \\
7 & 6 & Inverse eigenvalues and the characteristic-polynomial relation
for a contragredient, giving the formal reciprocal functional equation.
This calculation alone contains no absolute-value estimate. \\
9 & 8 & Passage to a motive over a finite field.  The determinant-weight
calculation explicitly assumes a motive satisfying the Weil conjectures;
the later calculation also invokes a dual/Tate-twist hypothesis. \\
\bottomrule
\end{tabular}
\end{center}

For scan 9, the visible words \emph{satisfaisant aux conj. de Weil} are
decisive to the scope.  The page uses geometric weight information when
passing from determinant algebra to a weight-normalized functional
equation.  It does not derive that information from an arithmetic
critical-line theorem.  Scan 4, archival page 3, is an upside-down reused
printed verso and was inspected only to understand the foliation; it is
not cited as an additional mathematical result.

For precision, the following finite-dimensional calculation states
exactly what the determinant mechanism can and cannot supply.  It is an
elementary consequence proved here, not attributed as a newly discovered
theorem in the manuscript.

\begin{proposition}[AR01: reciprocal symmetry does not imply purity]
Let \(a>0\) and \(c>0\).  On a two-dimensional real vector space let
\(f\) have eigenvalues \(c\) and \(a/c\).  Then
\[
 P_f(t)=\det(1-tf)
       =1-(c+a/c)t+at^2
       =at^2P_f\bigl((at)^{-1}\bigr).
\]
The two eigenvalues both have absolute value \(\sqrt a\) if and only if
\(c=\sqrt a\).  In particular, the determinant functional equation,
even with positive real eigenvalues and fixed determinant \(a\), does
not imply that absolute-value assertion.
\end{proposition}

\begin{proof}
The determinant is the product \((1-ct)(1-(a/c)t)\).  Expanding this
product gives the displayed quadratic; substituting \((at)^{-1}\)
and multiplying by \(at^2\) gives the same quadratic.  Because the
eigenvalues are positive, their absolute values equal themselves.
Both equal \(\sqrt a\) precisely when \(c=\sqrt a\), which also
forces \(a/c=\sqrt a\).  Choosing any other positive \(c\) proves
the final assertion.
\end{proof}

This counterexample concerns the determinant identities alone.  It does
not assert that the displayed linear map is the Frobenius of a pure
motive.  The point is to identify the missing hypothesis accurately:
purity is additional geometric input, exactly as scan 9 indicates.
Consequently these concrete archive hits clarify Q01--Q02 and the
\(\lambda\)-operation discussion without supplying a new transfer from
assumption \(A\) to arbitrary Hasse--Weil functions, regulators, or
integral \(K\)-groups.

\paragraph{Limits of this archive check.}
Cotes 15 and 70 were not inspected in this run.  The other supplied
locators, including cote 119 scan 65 and cote 162-1 scan 4, remain
text-record references here; no direct scan verification is claimed for
them.  The SA20 intake's cotes 119--121 and 153/155 do not by themselves
give an exact scan for an individual \(K\)-group or special-value
assertion.  No such locator has been invented.  Source-version differences
in Deligne--Goncharov, corrections to intake bounds, and incomplete Lean
mapping are not Grothendieck transcription corrections.


\section{Checked formal-group consequences and recovered source statements}
The new proofs ALP01 and ALP02 are checked below. ALP02 uses the explicitly stated intermediate source theorem, whose local proof and definitions were read; it is not deduced from the endpoint assumption A, and no new kernel certification is claimed.

\subsection{An unconditional algebraic intermediate}

The following standard derivation identity isolates an independently usable
part of the invariant-derivative mechanism.  The proof below is complete;
it does not require RH, a CM hypothesis, or a Lean build.

\begin{proposition}[ALP01: prime-iterate divisibility at an augmentation]
Let \(A,R\) be commutative rings, let \(\epsilon:A\to R\) be a ring
homomorphism, and let \(D:A\to A\) be an additive derivation.  If \(p\)
is prime and \(\epsilon(F)=\epsilon(G)=0\), then
\[
 \epsilon\bigl(D^p(FG)\bigr)\in pR.
\]
If \(A\) has characteristic \(p\), then \(D^p\) itself is a derivation.
\end{proposition}

\begin{proof}
Additivity and the Leibniz rule imply, by induction on \(n\),
\[
 D^n(FG)=\sum_{j=0}^{n}\binom nj D^j(F)D^{n-j}(G).
\]
For the induction step apply \(D\), use additivity and the Leibniz rule on
each summand, and combine adjacent coefficients using Pascal's identity.
An additive map commutes with multiplication by integers, so this argument
works over arbitrary commutative rings.  When \(n=p\), the terms with
\(j=0,p\) have zero image under \(\epsilon\), while
\(p\mid\binom pj\) for every \(0<j<p\).  Applying \(\epsilon\)
therefore puts every remaining summand in \(pR\).  In characteristic
\(p\) the intermediate terms vanish already in \(A\), leaving
\(D^p(FG)=F D^p(G)+D^p(F)G\).  Iteration preserves additivity, proving
the final assertion.
\end{proof}

The official source contains the constant-coefficient specialization as
\path{GlobalTheory.LawDer.constantCoeff_iterate_invDer_mul}, line 38
of
\path{Development/AbelianVarietyTheory/Proofs/DivisorProofs/ConstantCoeffIteratePInvariantDerivationMul.lean}.
Its hypotheses give an integral formal law with zero constant terms and a
right identity, the associated invariant derivation, and two power series
with zero constant terms.  The entire short local proof was read: it uses
the same iterated Leibniz formula and prime binomial divisibility.  This
establishes the stated algebraic mechanism; it does not by itself establish
\(D^p=0\) on a CM curve or a repeated theta-derivative valuation estimate.

\subsection{Targeted recovery from the official Lean archive}
\label{sec:alpoge-official}

The initial attachment-only extraction was followed by a second approach.
The hyperlink annotations in the supplied note, p.~1 footnote 1, point to
\url{https://alpo.ge/lean.tar.gz.part01} through
\url{https://alpo.ge/lean.tar.gz.part14}.  Ordinary GET requests retrieved
these public parts, totalling 239,275,554 bytes.  Streaming the concatenated
gzip archive counted 175,187 Lean files, matching the supplied aggregate.
Only targeted source files were materialized.  The concatenation hash,
per-part URLs and hashes, and selected-file hashes are recorded in the
retrieval manifest.  The complete archive is not bundled with this report.
No archive code, \path{lean}, or \path{lake} was executed.

This replaces the earlier absence-of-source gate for the definitions
actually recovered below.  It does not replace a statement-level reading
by a kernel certificate, or certify the imported proof chains.  The final
RH/Dirichlet results remain the stipulated theorems \(A\).

\paragraph{The CM-curve predicate has a different prime parameter.}
In
\path{Development/AbelianVarietyTheory/PublishedTheorems/}\texttt{Divisors\textsubscript{2}}\path{/CmCurveGoodOutsidePrimePremises.lean},
\path{IsGoodCMCurve} is defined at line 31 and
\path{CMCurveGoodOutsidePrimeSt} at line 35.  For every prime
\(p>3\) with \(p\equiv3\pmod4\), the latter asks for a Weierstrass
curve over its \(j\)-field, with CM by the integer lattice of
\(\mathbb Q(\sqrt{-p})\), good reduction outside \(p\), and minimal
discriminant \(-p^3\) up to local units after admissible changes of variables.
The declaration
\path{GlobalTheory.CMModel.cmCurveGoodOutsidePrime}, line 47 of
\path{Development/AbelianVarietyTheory/Proofs/DivisorProofs/CmCurveGoodOutsidePrime.lean},
has precisely that type.  Here \(p\) is ramified in the defining quadratic
field.  It is not the inert prime in a fixed CM field appearing in the
note's step 2.  Thus these candidate names cannot substitute for the
theta-derivative statement.

\paragraph{A genuine explicit determinant-depth theorem.}
A more relevant source statement was located and its hypotheses expanded.
Let \(C\) be an algebraically closed field of characteristic zero, let
\(v:C\to\mathbb R\cup\{+\infty\}\) be an additive valuation, and
let \(p\geq3\) be a prime with \(v(p)=1\).  Let \(d\in\mathbb N\),
let \(h\in\mathbb N\) be positive and even, and let
\(\mathrm{ty}:\{1,\ldots,d\}\to\mathbb Z/h\mathbb Z\).
Parity of a type means parity of its representative in
\(\{0,\ldots,h-1\}\); because \(h\) is even, this is well defined.

The predicate \path{LogFGL v d} consists of power series
\(\log_G=(\log_{G,a})\) with zero constant term and identity linear term,
and an integral formal law \(F\) with zero constant term satisfying
\(\log_G(F(x,y))=\log_G(x)+\log_G(y)\).
The input \path{MuAction G p h ty} gives formal maps \(M_\zeta\) with
zero constant term, integral coefficients whenever
\(\zeta^{p^h-1}=1\), and
\[
 \log_{G,a}(M_\zeta(x))
 =\zeta^{p^{\mathrm{ty}(a)}}\log_{G,a}(x)
 \qquad(\zeta^{p^h-1}=1).
\]
These exact structures occur at lines 37 and 89 of
\path{ReciprocityTheory/DualLattices/VEqZeroOfPowEqOne.lean}.

Choose \(j\in\mathbb N\) such that each odd type occurs at most \(j\)
times.  Choose an even type \(\tau\), an integer \(i\in\mathbb N\),
an integer \(c\geq1\), and a matrix \(u=(u_{ae})\in C^{i\times d}\).
Assume all entries of \(u\) are integral, one \(i\times i\) minor has
valuation zero, and \(u_{ae}=0\) whenever \(\mathrm{ty}(e)\ne\tau\).
The integrality and unit-minor condition is exactly
\path{IsSatFrame v u}.

The source defines \(\operatorname{detDepth}(G,u,c)\) as the infimum,
over integral formal power series on \(i\) copies of \(G\), of the
valuation obtained by composing with formal exponentials in the frame
directions and applying the ordinary differential operator
\(\det(\partial/\partial t_{ba})^c\) at zero.  Its definition includes
the multi-index factorials of ordinary differentiation, not divided Taylor
coefficients.  The exact definitions are \path{derivAt0},
\path{depth}, \path{detPow}, \path{IsSatFrame}, and
\path{detDepth} at lines 16, 36, 41, 45, and 49 of
\path{Development/AbelianVarietyTheory/Definitions/Foundations/IsSatFrame.lean}.

The cited source theorem asserts that there is an integer \(r_0\) with
\[
 0\leq r_0\leq i,
 \qquad
 r_0\leq j\left(\frac32+\frac12\log_p c\right),
\]
such that
\[
 \boxed{
 \operatorname{detDepth}(G,u,c)
 \geq
 \frac{(i-r_0)c p}{p^2-1}
 -i^2(\log_p c+1).}
\]
Here \(\log_p c=\log(c)/\log(p)\); \(c\) is a positive integer,
\(i\) may be zero, and the depth takes values in the extended real line.
All hypotheses just listed are present in the declared theorem.

The exact exported declaration is
\path{GlobalTheory.T.detDepth_lower_bound_saturated_even_type},
line 8 of
\path{Development/AbelianVarietyTheory/Proofs/DivisorProofs/DetDepthLowerBoundSaturatedEvenTypeStatement.lean}.
The local proof is
\path{GlobalTheory.DepthLow.detDepth_lower}, line 199 of the corresponding
\path{DetDepthLowerBoundSaturatedEvenType.lean}.
\emph{Check status:} the complete statement, these definitions, and the
entire local proof were read; the imported typical-curve, basis-change,
radius and divided-power-content lemmas were not independently re-proved,
and no kernel run is claimed.  This is a cited source theorem with an
explicit numerical bound, not a theorem inferred from \(A\).

\begin{corollary}[ALP02: rank-one specialization of the cited depth theorem]
Under the preceding hypotheses, suppose no odd type occurs, take \(i=1\),
and let \(u\in C^d\) be integral with a unit coordinate and supported in
the even type \(\tau\).  For every integral formal series \(f\) on
\(G\) and every integer \(c\geq1\), the cited source theorem gives
\[
 v\left(\left.\frac{d^c}{dt^c}f(\exp_G(tu))\right|_{t=0}\right)
 \geq \frac{cp}{p^2-1}-\log_p c-1.
\]
\end{corollary}

\begin{proof}
Set \(j=0\).  The asserted bound on the nonnegative integer \(r_0\)
forces \(r_0=0\).  Substituting \(i=1\) gives the right-hand side
displayed above as a lower bound on \(\operatorname{detDepth}(G,u,c)\).
For \(i=1\), the determinant polynomial is the \(c\)-th power of the
single differential variable.  By the definition of depth as an infimum,
each integral series has derivative valuation at least this lower bound.
\end{proof}

ALP02 is a formal consequence of the cited source theorem with its full
hypotheses retained.  It exposes the leading
\(p/(p^2-1)=1/p+1/p^3+\cdots\) and an explicit logarithmic error.
To identify its expression with \(D^c\theta(0)/\Omega^c\), one must
still supply the CM integral formal model, the indicated root-of-unity
action, the integral theta trivialization, a saturated tangent vector,
and the comparison of that vector with the chosen period \(\Omega\).
Those identifications are not supplied by the predicate
\path{CMCurveGoodOutsidePrimeSt}, and are not silently assumed here.

\paragraph{The actual Bessel-period predicate.}
\path{Formalism/PeriodIntegrals/BesselPeriodGGP.lean}, line 105, defines
\path{PI.BesselPeriodGGP.BesselPeriodGGP}.  It quantifies over an even
integer \(k\geq2\), \(r\geq0\), the indicated nondegenerate geometric
data \(z,c\), Harish--Chandra parameters \(u\) and
\(\lambda=(v,-b)\), the explicit strict interlacing and half-integral
difference conditions, a complete Hilbert discrete-series representation
\(\sigma\), a finite-dimensional representation \(\pi\) of the
specified compact type, a lowest unit vector, and Haar measures on the
specified unitary and unipotent subgroups.  Its conclusion is the
existence of two translates for which an operator-valued Bessel integral
is nonzero.  It asserts neither a literal value \(1\) nor a universal
vanishing alternative for non-interlacing parameters.  In particular,
the adjacent \path{GapBig} definition is not an additional argument in
this particular predicate.

The declared equivalence
\path{besselPeriod_iff}, at line 382 of
\path{Development/AbelianVarietyTheory/Proofs/JetHarmonics/SymmetricMeasureL2NormEqResolutionL2.lean},
identifies it with the separately defined
\path{FormalDegreeHaarQuotient.BesselPeriodNeZero}.  This is a useful
normalization/typing comparison, not itself a proof of either predicate.
The source-level definition is now recovered; a general global period
formula with values \(0,1\) has not been extracted.

\paragraph{The actual theta-lift interface.}
\path{UpperTheory.TP.TLiftK} is defined at line 12 of
\path{Development/AutomorphicArithmetic/Foundations/GaloisRepresentations/ThetaParseval/HalfIndexLattices/TLiftK.lean}.
It is an eventual-rank lift-existence statement over a CM field: for the
specified good base and a definite Hermitian space of dimension \(n\geq1\),
there is a threshold \(d_0\); for all \(d\geq d_0\), suitable definite
skew-Hermitian data, Hecke characters with the stated restriction and
infinity-type conditions, and an automorphic representation whose weights
satisfy \(d+m_W(w)\leq2\lambda_w(k)\), it produces an automorphic
subspace satisfying \path{IsLift}.  The threshold is existential.
The witness \path{UpperTheory.TP.K.tLiftK} is at line 4 of the
\path{ThetaParseval/ThetaKernels/TLiftK.lean} file under the same
\path{Development/AutomorphicArithmetic/Foundations/GaloisRepresentations/}
prefix.  This corrects the abbreviated path in the supplied map.
The source also exposes the \path{PropsB} fields in
\path{TypeConditions/WeightsConjBis.lean}: continuity, semisimplicity,
specified local compatibility, Hodge--Tate weights and conjugate weights,
and qualified crystalline/de Rham conditions.  None is a numerical
assertion that an arbitrary chain of \(m\) lifts gives multiplicity \(m\)
with a uniform cost.  That translation remains distinct from the recovered
lift-existence interface.



\clearpage
\part{Unchecked intermediate claims and the remaining gates}
\noindent This part inventories the source claims whose complete hypotheses,
proofs, or formal interfaces were not available for a fresh certification.
Their documentary locators are checked where stated. They are not promoted
to new independent theorems by acceptance of a final zero-location result.
The full statements and statuses are in \path{METHOD_STATEMENTS.tsv}.


\section{What the two routes supply: the Alp\"oge evidence}
\label{sec:alpoge-evidence}

\subsection{Theorems stipulated, declarations inspected, and proofs actually checked}
The standing hypothesis \(A\) is the user's instruction to cite Alp\"oge's
final RH and Dirichlet-GRH results as theorems.  This convention is separate
from a claim that the present run executed Lean.  We inspected every line of
\path{ALPOGE_Main.lean} and \path{ALPOGE_Statements.lean}.  The former has
five declarations: \path{riemannHypothesis},
\path{riemannZeta_eq_zero_re},
\path{LFunction_ne_zero_of_half_lt_re},
\path{LFunction_eq_zero_re}, and
\path{LFunction_eq_zero_re_of_isPrimitive}.  The last three respectively
state nonvanishing to the right of \(1/2\) (apart from the principal-character
pole convention), zero location in \(0<\Re s<1\), and its primitive-character
version.  In particular, they are statements about Dirichlet characters over
\(\mathbb Q\); an occurrence of general automorphic or CM-field terminology in
an internal construction does not enlarge their conclusion to all Hecke
characters.

The attached scan is a small aggregate table.  It reports zero occurrences of
\path{sorry}, \path{admit}, explicit axiom declarations and
\path{native_decide}; it is not a kernel-verification log.  The attached
\path{ALPOGE_LEAN_MAP_20261120.md}, opening paragraph and \S\S1--2, explicitly
says that its author performed text parsing only, never ran \path{lake} or
\path{lean}, and estimated the declaration dependency closure.  The actual
175,187-file archive was not among these attachments.  The official linked
multipart archive was subsequently retrieved and targeted source definitions
and proofs were read, as recorded in the official-source addendum.  We
record the final results as stipulated cited theorems, distinguish the map
from directly recovered source, and record no check event as a new formal
certificate.

\subsection{Statement-level inventory}
\label{sec:alpoge-table}
The accompanying \path{METHOD_STATEMENTS.tsv} preserves the complete table,
including the exact reported file paths.  The following distinctions are
mathematically material.

\begin{enumerate}
\item \textbf{Step 2, note p.~3.} The claim concerns an invariant derivation
\(D\) on a CM curve at an inert prime, \(D^p\equiv0\), and a valuation bound
for \(D^a\theta(0)/\Omega^a\) with leading exponent
\(a/p+a/p^3+\cdots\), up to \(O(\log a)\).
The map, \S4 row 2, supplies only candidate folders
\path{Development/AbelianVarietyTheory/BlockTheory/HeightDefect},
\path{Proofs/HomogeneousIdeals}, and \path{Proofs/CubicalLineBundles}, and
predicate names \path{IsGoodCMCurve}, \path{CMCurveGoodOutsidePrimeSt}.
It marks the match ``low''.  Direct recovery shows that those two predicates
concern good reduction outside the ramified prime defining
\(\mathbb Q(\sqrt{-p})\), a different role of \(p\).  The official-source
addendum recovers instead an explicit formal-group determinant-depth lower
bound and its rank-one derivative consequence, with an explicit logarithmic
error.  Its comparison with \(D^a\theta(0)/\Omega^a\) still requires the
CM formal model, theta trivialization and tangent/period normalization.
No theorem for arbitrary crystals or CM motives is inferred from that bound.

\item \textbf{Steps 3 and 5, note pp.~3--4 and p.~3 footnote 3.}
The path-quadric identity, definite middle-cohomology sign, iterated Parseval
estimate and prohibition of cancellation are described informally.  The map,
\S\S4--5, does not identify an unconditional theorem with this combined
statement.  Its quoted \path{LateralTheory.EWeil.Goal}, in
\path{Development/EtaleCohomology/WeilConjectures/FermatCount.lean:75}, has
an explicit antecedent \path{Base}; the same section reports no declaration
with \path{Goal} as its concluded type in its text search.  This is evidence
about the supplied map, not a proof that no such theorem exists in the full
archive.  We use the established Weil I/II theorems for finite-field purity and
do not count a new proof of them as extracted here.

\item \textbf{Steps 4 and 7, note pp.~3--4.}
The positivity assertions are \(W(1_S)\geq0\) and determinant/period
expressions involving interlacing weights.  For step 4 the map finds no named
bootstrap declaration.  For step 7 it gives
\path{Formalism/PeriodIntegrals/BesselPeriodGGP.lean} and
\path{Formalism/SharpRestriction/CuspSpanNormBounds/OscillatoryIntegralBounds.lean},
without a complete statement or normalization of the period.  Direct source
reading now identifies \path{BesselPeriodGGP} as a predicate asserting
nonzero operator-valued Bessel integrals for some translates under explicit
interlacing, type, representation and Haar hypotheses; the official-source
addendum records these.  Thus no universal assertion that a global
Gan--Gross--Prasad period is literally
\(0\) or \(1\) is extracted.  The classical compact branching multiplicity
\[
 \dim\operatorname{Hom}_{U(n-1)}(V_c,V_a|_{U(n-1)})
 =\begin{cases}1,&a_1\geq c_1\geq a_2\geq\cdots\geq c_{n-1}\geq a_n,\\
 0,&\text{otherwise},\end{cases}
\]
is already unconditional; see A.~W. Knapp, \emph{Branching theorems for compact
symmetric spaces}, Representation Theory \textbf{5} (2001), \S1, equation
(1.1), p.~404, or \emph{Lie Groups Beyond an Introduction}, second edition,
Theorem~9.14.  A dimension of an intertwiner space is not a specified
adelic period integral.  Normalized compact-character integrals recover this
multiplicity, whereas scaling a vector in a nonzero global period scales the
period.  This elementary distinction prevents an unsupported normalization
claim.

\item \textbf{Step 8, note p.~4.}
The note says a chain of theta lifts between totally definite unitary groups
wins multiplicity \(m\).  The map, \S3 level 3, quotes an exact internal
implication
\[
 \mathtt{galoisRepresentations\_of\_thetaLift}
 :\ \mathtt{TLiftK}\longrightarrow\mathtt{Row0b\ recT},
\]
at
\path{Development/AutomorphicArithmetic/Foundations/GaloisRepresentations/IntegralStructures/GaloisRepresentationsOfThetaLift.lean:6}.
It reports a term \path{UpperTheory.TP.K.tLiftK : TLiftK} at
\path{Development/AutomorphicArithmetic/Foundations/GaloisRepresentations/ThetaParseval/ThetaKernels/TLiftK.lean:4}
and \path{tLift_holds_k} at \path{TPGHolds.lean:58}.
The displayed predicate \path{Row0b}, in
\path{Foundations/GaloisRepresentations/TypeConditions/WeightsConjBis.lean:69},
requires a CM field, a regular algebraic cuspidal representation of
\(\mathrm{GL}_a\), a polarized pair, and concludes existence of an
\(\ell\)-adic representation satisfying \path{PropsB}.  Direct source
recovery now supplies both \path{PropsB} and \path{TLiftK}; the latter
is a lift-existence statement with an existential rank threshold and explicit
weight and character conditions.  It is not the claimed chain-multiplicity
estimate with all constants.  The addendum gives the precise distinction.

\item \textbf{Steps 9--10 and 12--16, note p.~4.}
The map identifies names for mass formulas, divisibility, seesaw/Rallis,
slopes, covolumes and orbital bounds.  The exact internal final interface it
quotes, in
\path{ThetaCorrespondence/ThetaPairLifts/ThirdPlaceSeries/SiegelEisensteinSecondTermHypotheses.lean:69},
requires seven inputs: \path{UnitaryGroupThirdPlaceHypotheses},
\path{ThetaLocalInputs}, \path{ThetaIntegralsAgree},
\path{LocalMatching}, \path{CovolumeBoundOfX}, \path{RecTExists},
and \path{VWin}; it concludes \path{EveryCharacter}.
The displayed \path{UnitaryGroupThirdPlaceHypotheses} itself has fields
\path{ClozelCohomological}, \path{GreenCuspidal}, and
\path{HeckeLAdicCharacter}.  The map reports internal discharge of these
interfaces.  The downloaded archive makes targeted bodies available, but the
entire trace/mass/slope dependency chains have not been read or kernel
checked in this run.  Names alone do not supply missing hypotheses for
independently reusing a numerical inequality.
For step 13 the candidate is
\path{Formalism/LatticeSlopes/LatticeSlopesProofs/SlopeInequality.lean},
together with \path{Formalism/LatticeSlopes/Statement.lean}.
Bost's slope method, the Rallis formula and compact branching are established
background, not new conclusions of RH.  The note alone does not isolate an
additional general trace formula with a test-function class and convergence
conditions.

\item \textbf{Step 17 and footnote 6, note p.~4.}
The note mentions classes in
\(H^1(\mathcal O_E[1/\ell],\mathbb Q_\ell/\mathbb Z_\ell(-e))\), a
weakening of Vandiver, and possible integral \(K\)-group consequences.
The map, \S4 row 17 and \S8, has no exact Vandiver/Lichtenbaum declaration and
only candidate folders \path{GroundTheory/ProfiniteCohomology/RestrictedRamification}
and \path{Development/AutomorphicArithmetic/CyclotomicBounds/ClassNumberBounds}.
No integral \(K\)-group vanishing or cyclicity assertion is promoted to a
checked theorem here.  This does not question the stipulated final RH result.

\item \textbf{Step 20, note p.~4.}
The exact public endpoint is the Dirichlet conclusion already used in the
main paper.  The map \S3 quotes
\path{UpperTheory.FinalTheorems.everyCharacter} at
\path{Development/AutomorphicArithmetic/CharacterBoxes/MainConclusion.lean:10}
and \path{allChi}, \path{rhOfAll}, \path{critAll},
\path{claimAll} in \path{Supplements/Kernels/MainTheorems.lean:7--16}.
Those statements match the actual attached \path{Main.lean} interfaces.
\end{enumerate}

\paragraph{Reuse independently of RH.}
An intermediate theorem with its own proof can be used without first invoking
RH, even when it occurs in a proof of RH.  The converse is false: accepting a
final theorem does not determine which lemmas, definitions, normalizations or
constants a proposed proof established.  Here compact branching and the
Deligne--Goncharov results below are independently usable cited background;
the official-source follow-up additionally supplies a precise determinant-depth
statement, its explicitly qualified rank-one consequence, and exact Bessel
and theta-lift predicates.  The remaining gates are their theta/period
normalizations and quantitative chain interpretation, rather than an absence
of all source.  Neither an internal constant
\(401\) in \path{VWin} nor the bound \(0<\varepsilon<1/52\) quoted in the map
is an effective class-number constant until its surrounding predicates have
been interpreted.



\section{Masullo: unchecked structure and trace claims}

The following exact locators identify candidate tools separately from RH.
They are recorded as draft claims; accepting \textbf{(M)} alone does not
prove the intermediate functional-analytic assertions.
{\small
\begin{longtable}{>{\raggedright\arraybackslash}p{0.19\linewidth}>{\raggedright\arraybackslash}p{0.21\linewidth}>{\raggedright\arraybackslash}p{0.50\linewidth}}
\hline
Locator & Tool & Exact restriction or status \\
\hline
\endhead
2.2.3.2, p.~82 & Adelic Riemann--Roch & Fourier/Poisson identity for adelic
theta functions; classical Tate-theory input, not an RH consequence.\\
3.2.1.3, p.~123; 3.2.2.1, pp.~124--126 & Hecke continuation and idelic Mellin
transform & Classical analytic continuation, functional equation, and
theta--zeta relation; the draft cites Tate theory.\\
5.2.1.1, p.~230; 5.2.3.1, p.~233 & First structure theorem & Strict LF
topology for theta spaces and nuclear Fr\'echet isotypes.\\
5.4.1.1, p.~265; 5.5.3.2, p.~282; 5.5.3.4, p.~285 & Second and third
structure theorems & S-local description and dual eigendistributions,
with the source's chosen sections and isotypes.\\
6.1.2.1, p.~288; 6.2.1.1, p.~295 & Fourth and fifth structure theorems &
S-local-to-global comparisons; the hypotheses include
$\operatorname{Pic}(\mathcal O_{F,S})=1$.\\
7.3.1.6, pp.~344--345 & Orbital Schauder bases & Saturated $S$ and trivial
S-class group; unconditional dual bases and absolute biorthogonal bases
on the specified $p$-primary quotients.\\
7.3.3.4, p.~360 & Bases on Paley--Wiener levels & Specified S-local orbital
quotients and support conditions; not arbitrary cohomology spaces.\\
10.3.1.1, pp.~448--449 & Spectral nuclearity & Sobolev degree $s>1$;
trace multiplicity is truncated to
$\min(N,\lfloor(s-1)/2+1\rfloor)$.\\
11.2.2.1, p.~490 & Orbital nuclearity & $s\ge2[F:\mathbf Q]$ and
unramifiedness at all places above the chosen prime.\\
11.3.2.1, p.~495 & Orbital traces & Support in $C_F^{>1}$, prescribed
invariance and $\operatorname{Pic}(\mathcal O_{F,S})=1$.\\
12.4.1.1, p.~514 & S-local trace calculation & Truncation limit and local
principal-value terms, with source hypotheses on $S$ and invariance.\\
12.5.1.1, pp.~519--520 & Global trace formula & Requires
$\phi|_{C_F^1}=0$; uses a spectral strip and the chosen completions.\\
13.3.1.1, p.~522 & Proposed critical-line trace formula & Uses
$N_{\rm spec}=\min(N,[F:\mathbf Q])$ and Mellin derivative-vanishing
conditions; carries the source's correction footnote.\\
\hline
\end{longtable}
}

No comparison here constructs a Weil cohomology over $\operatorname{Spec}
\mathbf Z$, an algebraic correspondence realizing the analytic operator,
or the positivity in the standard conjectures. Such a realization would
be an additional bridge. In particular the restricted test-function
condition cannot simply be discarded.




\section{What remains after the reductions}
The mathematical questions of this run have explicit transfer answers,
refutations, or reduced gates. General non-CM motivic RH is not a theorem
of either stipulated input. For a given scheme the residual list in ZT06
also retains analytic continuation and functional equations wherever these
are not independently known. For the non-CM elliptic case these analytic
properties are known and ZT07 isolates RH itself without possible
off-line cancellation.

The outstanding source-extraction questions are different in kind. The
definition of $R$ behind SA06 A.8c is needed to identify its exact
class-number normalization; the maximal imaginary quadratic-order case
is already handled by IA02--IA03. The original proof behind LM8\_PUSH
A.69 is needed to determine its abbreviation DS/RH; its supplied
curve-genus context does not establish an analytic GRH premise. Original
source searches did not recover either full run. The recovered formal-group source theorem supplies the explicit derivative bound in ALP02. Its identification with a normalized CM theta value still requires the integral theta model, root-of-unity action, saturated tangent, and period comparison. The recovered Bessel and theta-lift predicates retain their exact nonvanishing and eventual-rank scope. Masullo's long analytic
structure and trace claims require their own proof verification; M alone
is not a certificate for them. These gates, all tried retrieval routes,
and the precise next actions are retained in the ledger.

The special-value gate is also distinct. A removes the quadratic
exceptional-zero condition needed in the CM-height argument, but does not
prove the remaining general individual Colmez identity. Exact Bernoulli
and regulator formulas, integral $K$-groups, and extensions remain their
own arithmetic statements. The counterexamples and computations in
Part I identify why they cannot be replaced by a zero-location slogan.

\clearpage
\begin{thebibliography}{99}
\addcontentsline{toc}{section}{Bibliography and exact source locators}
\bibitem{Alpoge} L. Alp\"oge, \emph{Remark}, supplied four-page version,
20 November 2026. Steps 1--20, pp.~3--4; actual endpoint declarations
\path{ALPOGE_Main.lean}, lines 4--16, with definitions in
\path{ALPOGE_Statements.lean}. Internal paths and their evidence status:
\path{ALPOGE_LEAN_MAP_20261120.md}, \S\S3--5 and 8.
\url{https://alpo.ge/remark.pdf}.
\bibitem{Masullo} A. M. Masullo, \emph{The Music of the Primes}, supplied
940-page draft, 20 November 2026. \S1.2, PDF p.~12; numbered claims and
pages in the method table; intended divisor theorem 13.3.1.2, PDF p.~522,
and \S14.1.2, PDF p.~528. The attached PDF is the version of record here.
\bibitem{Anderson1986} G. W. Anderson, \emph{Cyclotomy and an extension of
the Taniyama group}, Compositio Math. \textbf{57} (1986), 153--217.
\S2.1, p.~159; Proposition 3.8.1, p.~168; \S5.3, p.~178;
Corollaries 5.7.1--5.7.2 and Remark 5.7.3, p.~181.
\url{https://www.numdam.org/item/CM_1986__57_2_153_0.pdf}.
\bibitem{FarguesCM} L. Fargues, \emph{Motives and automorphic forms:
the (potentially) abelian case}, notes, \S7.1, Theorem 7.1, pp.~24--25;
\S7.2, Proposition 7.8 and Corollary 7.9, p.~26.
\url{https://webusers.imj-prg.fr/~laurent.fargues/Motifs_abeliens.pdf}.
\bibitem{Milne1990} J. S. Milne, \emph{Canonical models of (mixed)
Shimura varieties and automorphic vector bundles}, 1990,
Theorem 6.8 and Remark 6.9, p.~313.
\url{https://www.jmilne.org/math/articles/1990a.pdf}.
\bibitem{Deligne1979} P. Deligne, \emph{Valeurs de fonctions $L$ et
p\'eriodes d'int\'egrales}, Proc. Sympos. Pure Math. \textbf{33}, part 2
(1979), 313--346. \S\S1.1--1.3; (3.1.2); Conjectures 1.8 and 2.8
are special-value conjectures, not numbered RH assertions.
\bibitem{SGA4} M. Artin, A. Grothendieck, J.-L. Verdier et al.,
\emph{SGA 4}, tome 3, Lecture Notes in Math. 305 (1973),
Expos\'e XVII, \S\S5.1--5.3: proper-support direct image and base change.
\bibitem{SGA5} A. Grothendieck et al., \emph{SGA 5}, Lecture Notes in
Math. 589 (1977), Expos\'e XV, \S3, n\textsuperscript{o}~2,
unnumbered Th\'eor\`eme, equation (1); n\textsuperscript{o}~3,
equations (2), (2$'$): rationality and trace expansion.
\bibitem{Bourbaki279} A. Grothendieck, \emph{Formule de Lefschetz et
rationalit\'e des fonctions $L$}, S\'eminaire Bourbaki, expos\'e 279
(1964/65), Theorem 5.1.
\url{https://www.numdam.org/item/SB_1964-1966__9__41_0.pdf}.
\bibitem{SGA6} P. Berthelot, A. Grothendieck, L. Illusie et al.,
\emph{SGA 6}, Lecture Notes in Math. 225 (1971), Expos\'e V,
\S2 and \S7, equation (7.1.1): lambda and Adams operations.
\bibitem{EGA4} A. Grothendieck and J. Dieudonn\'e, \emph{EGA IV$_3$},
Publ. Math. IHES \textbf{28} (1966), (8.8.2): finite-presentation descent.
\bibitem{WeilI} P. Deligne, \emph{La conjecture de Weil. I}, Publ. Math.
IHES \textbf{43} (1974), 273--307, Theorem 1.6.
\bibitem{WeilII} P. Deligne, \emph{La conjecture de Weil. II}, Publ. Math.
IHES \textbf{52} (1980), 137--252, Theorem 3.3.1.
\bibitem{Bittner} F. Bittner, \emph{The universal Euler characteristic
for varieties of characteristic zero}, Compositio Math. \textbf{140}
(2004), 1011--1032, Theorem 3.1 and its projective refinement.
\url{https://arxiv.org/abs/math/0111062}.
\bibitem{Silverman} J. H. Silverman, \emph{The Arithmetic of Elliptic
Curves}, second edition, Exercise VIII.8.19, especially part (b): local
Euler factors and the relation with the arithmetic surface zeta function.
\bibitem{LLS} Y. Lamzouri, X. Li, and K. Soundararajan,
\emph{Conditional bounds for the least quadratic non-residue and related
problems}, arXiv:1309.3595v3. Theorem 1.5, pp.~3--4; Corollaries 1.1--1.2,
pp.~2--3; corrected constants from this version.
\url{https://arxiv.org/abs/1309.3595}.
\bibitem{IMS} Y. Ihara, V. K. Murty, and M. Shimura,
\emph{On the logarithmic derivatives of Dirichlet $L$-functions at $s=1$},
Corollary 3.3.2, equation (3.3.3), p.~9.
\url{https://www.kurims.kyoto-u.ac.jp/~kenkyubu/emeritus/ihara/Publications-and-Recent-Preprints/RecentArticles/pdf-files/4A.pdf}.
\bibitem{HJM} P. Habegger, G. Jones, and D. Masser,
\emph{Six unlikely intersection problems in search of effectivity},
equation (5.9), Lemma 14(i), printed p.~22 / PDF p.~23.
\url{https://pure.manchester.ac.uk/ws/portalfiles/portal/36388863/final.logged.pdf}.
\bibitem{Wei} X. Wei, \emph{Smallness of Faltings heights of CM abelian
varieties}, arXiv:2111.00617v1. Definitions p.~2; Theorems 1.3--1.4,
pp.~4--5; Corollary 4.3, equation (13), p.~14; Proposition 4.7, p.~17;
proof of Theorem 1.3 and Theorem 4.8, p.~18; Proposition 5.3, pp.~19--20.
\url{https://arxiv.org/abs/2111.00617}.
\bibitem{DGsupersingular} C. Delfs and S. D. Galbraith,
\emph{Computing isogenies between supersingular elliptic curves over
$\mathbf F_p$}, arXiv:1310.7789, geometric count p.~1 and equation (1),
p.~2. \url{https://arxiv.org/abs/1310.7789}.
\bibitem{Miller} J. C. Miller, \emph{Real cyclotomic fields of prime
conductor and their class numbers}, arXiv:1407.2373, Theorem 1.1:
the $p\le151$ clause is unconditional.
\url{https://arxiv.org/abs/1407.2373}.
\bibitem{Washington} L. C. Washington, \emph{Introduction to Cyclotomic
Fields}, Theorem 4.17. Formula rechecked in D. S. Thakur,
\emph{Fermat's last theorem for regular primes}, in \emph{Cyclotomic
Fields}, printed p.~172 / PDF p.~179.
\url{https://www.sas.rochester.edu/mth/sites/dinesh-thakur/BPCyclotomic.pdf}.
\bibitem{Rim} D. S. Rim, \emph{Modules over finite groups}, Ann. of Math.
\textbf{69} (1959), 700--712, Theorem 6.24, p.~711. Original proof carried
as background; formulation and exact locator rechecked in S. Ullom,
\emph{Normal bases in Galois extensions of number fields}, p.~159,
and A. Siviero, \emph{Realisable classes, Stickelberger subgroup and its
behaviour under change of the base field}, Lemma 4.5, equation (16), p.~82.
\url{https://www.numdam.org/item/10.5802/pmb.13.pdf}.
\bibitem{Sutherland} A. V. Sutherland, \emph{18.785 Number Theory I},
Fall 2021, Lecture 19, Theorems 19.12 (p.~9) and 19.15 (p.~11).
\url{https://ocw.mit.edu/courses/18-785-number-theory-i-fall-2021/mit18_785f21_lec19.pdf}.
\bibitem{DLMF} NIST Digital Library of Mathematical Functions,
equations 25.15.3, 25.11.13, 25.11.14, 25.11.18.
\url{https://dlmf.nist.gov/25.15}; \url{https://dlmf.nist.gov/25.11}.
\bibitem{DeligneGoncharov} P. Deligne and A. B. Goncharov,
\emph{Groupes fondamentaux motiviques de Tate mixte}, supplied February
2003 TeX version: \path{Sec1ccc.tex}, \S1.6, Propositions 1.7 and 1.9;
\path{Sec2ccc.tex}, (2.1.3), Propositions 2.2--2.3, \S2.4,
\S2.17 and Proposition 2.18. Numbering is version specific.
\bibitem{Borel1974} A. Borel, \emph{Stable real cohomology of arithmetic
groups}, Ann. Sci. ENS \textbf{7} (1974), Proposition 12.2, p.~271.
\bibitem{Borel1977} A. Borel, \emph{Cohomologie de SLn et valeurs de
fonctions z\^eta aux points entiers}, Ann. Scuola Norm. Sup. Pisa
\textbf{4} (1977), Theorem 6.2 and \S6.4, pp.~631--633; 1980 erratum,
p.~373.
\bibitem{Knapp} A. W. Knapp, \emph{Branching theorems for compact
symmetric spaces}, Representation Theory \textbf{5} (2001), \S1,
equation (1.1), p.~404.
\url{https://www.math.stonybrook.edu/~aknapp/pdf-files/branching.pdf}.
\bibitem{Archive} Archives Grothendieck, Universit\'e de Montpellier:
cote 34, PDF scans 64 and 74 (manuscript sheets 63 and 73);
cote 60, \emph{Fonctions $L$ / \'Equation fonctionnelle des fonctions $L$
(cas g\'eom\'etrique)}, cover and scans 3, 5, 7, 9.
\url{https://grothendieck.umontpellier.fr/34.pdf};
\url{https://grothendieck.umontpellier.fr/60.pdf}.
\bibitem{PoonenTate2015}
B.~Poonen, \emph{Tate's thesis}, MIT 18.786 Number theory II lecture notes,
Spring 2015, \S4.6--4.7, pp.~18--19, and \S5.11,
Theorem~5.22, pp.~38--39.
\url{https://math.mit.edu/~poonen/786/notes.pdf}.

\bibitem{Shahidi1980}
F.~Shahidi, \emph{On nonvanishing of $L$-functions},
Bull. Amer. Math. Soc. (N.S.) \textbf{2} (1980), 462--464;
unnumbered theorem, p.~462.

\bibitem{Serre1972}
J.-P.~Serre, \emph{Propri\'et\'es galoisiennes des points d'ordre fini des
courbes elliptiques}, Invent. Math. \textbf{15} (1972), 259--331;
\S4.4, Th\'eor\`eme~3, p.~299.

\bibitem{SuppliedRuns} Supplied runs, 20 November 2026:
\path{BRIDGES_20261120.tex}, F1--F2 and Q01--Q02;
\path{INTAKE_RESEARCH_SA06_20261120.md}, A.8a, A.8c, B.4a--e, E.4a--p;
\path{INTAKE_RESEARCH_SA07_20261120.md}, the exact row mapping in
\path{CONDITIONAL_AUDIT.tsv};
\path{INTAKE_RESEARCH_SA20_20261120.md}, A5.011 and A.103;
\path{INTAKE_REVLOC_20261120.md}, I.5 and I.8b.
The complete supplied-source manifest is \path{audit/SOURCE_MANIFEST.json}.
\end{thebibliography}
\end{document}
